% gvim: encoding=utf8

\documentclass[11pt, french, screen, research-development]{report-rd-info}
   % - 11pt:  12pt peut être préférable pour faciliter la lecture sur les petits écrans, demander à l'encadrement
   % - french:  à remplacer par english en cas de rédaction (exceptionnelle) en anglais
   % - screen:  à enlever pour obtenir un rapport au format A4
   % - research-development:  à remplacer par 'research' seul, ou 'development' seul, voire 'intelligence' en fonction de la demande du jury à l'issue du travail
\usepackage[utf8]{inputenc}
\usepackage[T1]{fontenc}

% définitions propres au contenu actuel
\usepackage{amsmath, amssymb}
\usepackage{algorithm}
   \floatname{algorithm}{Algorithme}
   \ifenglish\renewcommand{\listalgorithmname}{List of Algorithms}
   \else     \renewcommand{\listalgorithmname}{Liste des algorithmes}
   \fi
\usepackage{listings}
   \lstset{backgroundcolor={\color{yellow}},
           columns=fullflexible,
           commentstyle={\color{blue}},
           numbers=left,
           numbersep=5pt,
           numberstyle={\tiny},
           stepnumber=1}
\usepackage{algorithmic}
   \renewcommand{\algorithmicrequire}{\textbf{Entrée~:}}
   \renewcommand{\algorithmicensure}{\textbf{Sortie~:}}
\usepackage{xcolor}
\usepackage{float}
\usepackage{url}
\usepackage{amsthm}
\usepackage{graphicx}
\usepackage{adjustbox}
   % - réduit les tableaux trop larges à la largeur de la colonne

\newtheorem{theoreme}{Théorème}
\newtheorem{preuve}{Preuve}
\newtheorem{definition}{Définition}

% Emplacement réservé à un contenu produit ultérieurement (schéma Excalidraw, image, résultat).
% Recherche de tous les emplacements : grep -rn "placeholder\|acompleter" documentation/rapport
\newcommand{\placeholder}[1]{%
   \fbox{\parbox{0.92\linewidth}{\centering\vspace{1.5em}\textcolor{red}{\textbf{[#1]}}\vspace{1.5em}}}}
% Information manquante, en ligne dans le texte ou dans un tableau.
\newcommand{\acompleter}[1]{\textcolor{red}{\textbf{[à compléter: #1]}}}

\begin{document}

\title{Graph Transformer pour la classification multi-label de molécules odorantes}
\subtitle{Rapport de phase~I}
\authorA{Tom}{Daniaud}
\authorB{Florian}{Petiot}
\supervisor{Fabrice}{Guillet}
\cosupervisor{Angélique}{Villière}
\coordinator{Prénom}{Nom} % [à compléter : coordinateur des PRED]
\institution{LS2N}
\theme{\'Equipe DUKE} % [à vérifier : équipe de rattachement de l'encadrant]
\coinstitution{Oniris Nantes}{Oniris}{1.7cm}
   % - nécessite le logo logos/Oniris.pdf (sinon commenter cette ligne)
\date{24 novembre 2026}

%--------------------------------------------------------------------------------

\begin{abstract}
\small
L'identification des odeurs des molécules constituant un arôme repose aujourd'hui sur l'olfactométrie, qui combine analyse chimique et perception d'un juge entraîné.
Prédire ces odeurs à partir de la seule structure moléculaire est un problème de classification multi-label difficile~: les données sont rares (environ 5\,000 molécules), bruitées et très déséquilibrées.
Les modèles de référence (POM, OpenPOM) reposent sur des réseaux de neurones sur graphes 2D, myopes vis-à-vis des atomes éloignés et aveugles à la géométrie.
Nous étudions l'apport d'un transformeur de graphes 3D pré-entraîné, Uni-Mol, combiné à un décodeur multi-label de type MLP, sur le jeu de données GS-LF (4\,983 molécules, 138 descripteurs).
Nous avons construit une chaîne complète~: génération de conformères 3D, encodage au format Uni-Mol, augmentation géométrique, perte asymétrique pour le déséquilibre, dégel partiel du transformeur, et un protocole d'évaluation aligné sur celui de la littérature.
Dix-sept expérimentations successives ont porté le F1 macro de validation de 0,21 à 0,35 et l'AUROC macro jusqu'à 0,87, un ordre de grandeur proche des modèles de référence, mais pas encore directement comparable faute d'évaluation sur un test tenu à l'écart.
La phase~II portera sur l'évaluation selon ce protocole, le passage à Uni-Mol2, la prise en compte des dépendances entre descripteurs et l'agrégation de nouvelles données, notamment stéréochimiques.
\end{abstract}

\begin{classification}
\small
   \category{I.2.6}{Artificial Intelligence}{Learning -- Connectionism and neural nets}
   \category{I.5.2}{Pattern Recognition}{Design Methodology -- Classifier design and evaluation}
   \category{J.2}{Physical Sciences and Engineering}{Chemistry}
   \terms{Algorithms, Experimentation, Performance, Measurement.}
   \keywords{Classification multi-label, olfaction computationnelle, relations structure-odeur, transformeur de graphes, Uni-Mol, représentation moléculaire 3D, apprentissage par transfert, déséquilibre de classes.}
\end{classification}

\maketitle

%--------------------------------------------------------------------------------

\begin{acknowledgements}
   Nous remercions Fabrice Guillet (LS2N) et Angélique Villière (Oniris) pour leur encadrement, leur disponibilité et les échanges qui ont orienté ce travail.
\end{acknowledgements}

%--------------------------------------------------------------------------------

\newpage

\tableofcontents

%--------------------------------------------------------------------------------

\chapter{Introduction}
\label{chap:Introduction}

Ce chapitre présente la problématique de la prédiction des odeurs à partir de la structure des molécules, les objectifs que nous nous sommes fixés, le travail réalisé durant la phase~I et ses principaux apports.
Il se termine par le plan du rapport.

%--------------------------------------------------------------------------------

\section{Présentation de la problématique}
\label{sec:Problematique}

L'équipe <<~Flaveur~>> du laboratoire GPEA d'Oniris étudie les arômes des aliments par olfactométrie.
Cette technique couple une séparation chimique des composés volatils à la perception humaine d'un juge entraîné, un <<~nez~>>, qui décrit l'odeur de chaque molécule au moment où elle est détectée \cite{Meierhenrich-2005}.
L'identification est donc coûteuse, lente et dépendante de l'expertise humaine.
Disposer d'un modèle capable de proposer les descripteurs olfactifs d'une molécule à partir de sa seule structure permettrait d'assister ce travail.

Ce problème relève des relations structure-odeur, étudiées depuis longtemps en chimie, et que l'apprentissage automatique aborde aujourd'hui comme un problème de \emph{classification multi-label}~: une molécule est associée simultanément à plusieurs descripteurs (\emph{fruité}, \emph{vert}, \emph{boisé}, etc.) choisis dans un vocabulaire fixe.
Deux difficultés le rendent particulièrement délicat~:
\begin{itemize}
   \item \textbf{la rareté des données}~: les jeux de données annotés publics comptent quelques milliers de molécules seulement (4\,983 pour le jeu GS-LF utilisé ici)~;
   \item \textbf{leur variabilité}~: la perception olfactive varie d'un individu à l'autre (physiologie, anosmies, culture, représentation mentale des odeurs), ce qui rend les annotations bruitées~; s'y ajoute un fort déséquilibre, certains descripteurs étant attribués à près de 2\,000 molécules et d'autres à une trentaine.
\end{itemize}

Les approches de référence \cite{SanchezLengeling-2019, Lee-2023} représentent la molécule par son graphe (atomes et liaisons) et l'encodent avec un réseau de neurones sur graphes (GNN).
Ces réseaux procèdent par propagation de messages entre atomes voisins~: ils sont naturellement myopes vis-à-vis des atomes éloignés dans le graphe et ignorent la forme tridimensionnelle de la molécule, alors que cette dernière conditionne son interaction avec les récepteurs olfactifs.
Le mécanisme d'attention des transformeurs \cite{Vaswani-2017} ne souffre pas de cette limite de distance, et les transformeurs de graphes \cite{Shehzad-2024, Kim-2022} l'adaptent aux données non séquentielles.
Parmi eux, Uni-Mol \cite{Zhou-2023} et Uni-Mol2 \cite{Ji-2024} sont pré-entraînés sur des centaines de millions de conformations 3D et intègrent explicitement la géométrie.
La problématique de ce travail est donc la suivante~:
\begin{quote}
   \emph{un transformeur de graphes 3D pré-entraîné, adapté à la classification multi-label, permet-il de mieux prédire les descripteurs olfactifs d'une molécule que les GNN 2D de référence, malgré la faible quantité de données disponibles~?}
\end{quote}

%--------------------------------------------------------------------------------

\section{Objectifs poursuivis}
\label{sec:Objectifs}

Le sujet fixe comme cadre l'utilisation du transformeur Uni-Mol, combiné à un décodeur multi-label de type MLP, sur le jeu de données GS-LF.
Nous en avons dérivé deux objectifs principaux.

\paragraph{Construire une méthodologie d'apprentissage.}
Il s'agit de mettre en place une chaîne complète et reproductible, de la molécule (code SMILES) jusqu'aux probabilités des 138 descripteurs~:
choix et préparation des données, transformation des molécules en entrées du transformeur, intégration d'un décodeur multi-label, choix du mode d'apprentissage (modèle gelé ou ajustement fin), traitement du déséquilibre des classes et choix des hyperparamètres.

\paragraph{Évaluer la précision du modèle.}
Il s'agit de mesurer les performances du modèle avec des métriques adaptées au multi-label déséquilibré (AUROC, AUPRC, F1, précision et rappel en moyenne macro) et, surtout, de les rendre \emph{directement comparables} à celles publiées pour les modèles de référence.
Cela impose de reproduire leur protocole d'évaluation, et non seulement leurs métriques.

%--------------------------------------------------------------------------------

\section{Travail réalisé}
\label{sec:TravailRealise}

Le travail de la phase~I s'est déroulé en trois temps.

\begin{enumerate}
   \item \textbf{Étude bibliographique}~: analyse des modèles de référence en olfaction (QSOR, POM, OpenPOM), du transformeur Uni-Mol et de sa seconde version, puis d'un travail récent appliquant Uni-Mol2 à l'olfaction \cite{Han-2026}, dont nous avons repris le protocole d'évaluation.
   \item \textbf{Développement}~: une chaîne d'apprentissage modulaire (PyTorch Lightning et Hydra), qui génère hors ligne des conformères 3D, les encode au format Uni-Mol, les augmente à la volée, et ajuste le transformeur avec un décodeur MLP.
   \item \textbf{Expérimentation}~: dix-sept expériences successives, chacune isolant un choix (fonction de perte, régularisation, ordonnanceur, augmentation, découpage, décodeur, dégel du transformeur), suivies dans TensorBoard et consignées dans un journal d'expériences.
\end{enumerate}

Une première version utilisait des représentations Uni-Mol figées calculées une fois pour toutes.
Elle a été remplacée par un encodage de conformères 3D, qui permet l'augmentation des données puis l'ajustement fin du transformeur.
Enfin, le protocole d'évaluation de la littérature a été implémenté en fin de phase~; son exécution complète reste à faire.

%--------------------------------------------------------------------------------

\section{Contribution}
\label{sec:Contribution}

Les apports de cette première phase sont les suivants~:
\begin{itemize}
   \item une chaîne complète d'apprentissage d'un classifieur multi-label d'odeurs fondé sur Uni-Mol, à partir de conformères 3D, avec augmentation géométrique, perte asymétrique \cite{Ridnik-2021} et dégel partiel du transformeur~;
   \item une progression mesurée du F1 macro de validation de 0,21 à 0,35 et de l'AUROC macro de 0,84 à 0,87 au fil des expériences, la plus forte amélioration venant de la correction de l'encodage des entrées d'Uni-Mol puis de l'ajustement fin de ses dernières couches~;
   \item la mise en évidence que nos premières mesures n'étaient pas comparables à la littérature (évaluation sur la validation, meilleure époque, seuil global), et l'implémentation du protocole de \cite{Han-2026}~: test tenu à l'écart, validation croisée à cinq plis, seuils par descripteur appris hors pli, ensemble des modèles et intervalles de confiance par bootstrap~;
   \item une analyse des limites de l'approche, en particulier l'invariance d'Uni-Mol par réflexion, qui l'empêche de distinguer deux énantiomères, et des pistes pour la phase~II.
\end{itemize}

À ce stade, nos meilleurs résultats sont du même ordre de grandeur que ceux d'OpenPOM (F1 macro de 0,375 sur le test), sans qu'une conclusion définitive puisse être tirée avant l'évaluation selon le protocole de référence.

%--------------------------------------------------------------------------------

\section{Plan de l'étude}
\label{sec:Plan}

Le chapitre~\ref{chap:EtatArt} étudie les travaux de la littérature~: les GNN de référence en olfaction (POM et OpenPOM), le transformeur Uni-Mol et sa seconde version, et l'application d'Uni-Mol2 à l'olfaction.
Leur analyse comparative fait ressortir les pistes que nous avons retenues.

Le chapitre~\ref{chap:Propositions} formalise notre proposition~: l'adaptation d'Uni-Mol à la classification multi-label des odeurs, le traitement du déséquilibre et le protocole d'évaluation.

Le chapitre~\ref{chap:Experimentations} décrit les outils et les développements réalisés, la mise en place des expériences, leurs résultats et leur analyse.

La conclusion résume le travail effectué, en tire les enseignements et présente les perspectives de la phase~II.
Les annexes regroupent les fiches de lecture, la planification, les fiches de suivi et l'auto-évaluation.

\chapter{\'Etat de l'art}
\label{chap:EtatArt}

Ce chapitre fait le point, à notre connaissance, sur les travaux utiles à la prédiction des descripteurs olfactifs à partir de la structure moléculaire.

Une molécule peut être représentée de plusieurs façons~: par une chaîne de caractères (code SMILES), par des descripteurs physico-chimiques ou des empreintes moléculaires (\emph{fingerprints}) calculés par des logiciels de chimie, par un graphe dont les sommets sont les atomes et les arêtes les liaisons, ou par la géométrie 3D de l'une de ses conformations.
Les modèles d'apprentissage profond récents opèrent principalement sur les deux dernières représentations.

Nous avons retenu trois familles de travaux~:
\begin{itemize}
   \item les GNN de référence en olfaction, QSOR, POM et sa réimplémentation ouverte OpenPOM, qui définissent l'état de l'art sur le jeu GS-LF et que le sujet désigne comme \emph{baselines} (section~\ref{sec:EAPom})~;
   \item le transformeur de graphes 3D Uni-Mol et sa seconde version Uni-Mol2, cœur du sujet (section~\ref{sec:EAUniMol})~;
   \item un travail récent qui applique Uni-Mol2 à l'olfaction sur le même jeu de données et dont le protocole d'évaluation sert de référence à nos mesures (section~\ref{sec:EAHan}).
\end{itemize}

D'autres transformeurs de graphes appliqués à l'olfaction ou aux propriétés moléculaires ont été identifiés, notamment GraphNOSE \cite{Sharma-2026} et une représentation à double graphe \cite{Zhang-2026}.
Leur étude détaillée est prévue en phase~II.
Les notions générales d'apprentissage profond (perceptron multicouche, rétropropagation, mécanisme d'attention \cite{Vaswani-2017}) sont supposées connues du lecteur.

%--------------------------------------------------------------------------------

\section{Les GNN de référence en olfaction~: POM et OpenPOM}
\label{sec:EAPom}

\subsection{Présentation}

\paragraph{QSOR.}
Sanchez-Lengeling et al. \cite{SanchezLengeling-2019} ont été les premiers à aborder la prédiction des odeurs comme une tâche de classification multi-label à grande échelle par apprentissage profond.
Ils assemblent environ 5\,000 molécules issues de deux catalogues de parfumerie, GoodScents et Leffingwell, annotées avec 138 descripteurs, et entraînent un GNN par propagation de messages directement sur le graphe moléculaire.
Le GNN surpasse les approches classiques fondées sur des empreintes moléculaires ou des descripteurs chimiques, et la représentation qu'il apprend (son \emph{embedding}) organise les molécules selon leurs odeurs~: elle se transfère à d'autres tâches olfactives.

\paragraph{POM.}
Lee et al. \cite{Lee-2023} prolongent ce travail avec la \emph{Principal Odor Map} (POM).
Le modèle est un GNN entraîné sur la fusion nettoyée de GoodScents et Leffingwell, appelée GS-LF~: c'est le jeu de données utilisé dans ce rapport.
La couche cachée du réseau définit une <<~carte des odeurs~>> dans laquelle des molécules aux odeurs proches sont proches.
Les auteurs valident le modèle de manière prospective~: sur environ 400 molécules nouvelles évaluées par un panel humain entraîné, ses prédictions sont, pour une majorité des molécules, plus proches du consensus du panel que celles d'un panéliste médian.
Le modèle est un ensemble de réseaux dont les hyperparamètres ont été choisis par une recherche de 500 essais en validation croisée.

\paragraph{OpenPOM.}
Le code et les poids de POM n'étant pas publiés, le groupe BioMachineLearning en a proposé une réimplémentation ouverte, OpenPOM\footnote{BioMachineLearning, \emph{OpenPOM}, \url{https://github.com/BioMachineLearning/openpom}}, fondée sur la bibliothèque DeepChem.
OpenPOM fournit le jeu GS-LF nettoyé (4\,983 molécules, 138 descripteurs), un réseau à propagation de messages et un découpage des données par stratification itérative multi-label d'ordre~2.
C'est aujourd'hui le point de comparaison reproductible de référence pour la tâche.
Ses performances sur le test de GS-LF, mesurées par \cite{Han-2026}, sont rapportées dans le tableau~\ref{tab:ResultatsLitterature}.

\begin{figure}
   \centering
   \placeholder{schema: GNN par propagation de messages -- molécule en graphe, agrégation des voisins sur K couches, lecture globale, MLP, 138 sorties~; mettre en évidence le voisinage limité à K liaisons}
   \caption{Principe d'un GNN par propagation de messages appliqué à une molécule}
   \label{fig:GNN}
\end{figure}

\subsection{Analyse}

\subsubsection{Intérêts des propositions}

Ces travaux établissent le cadre de notre problème~: la formulation multi-label sur 138 descripteurs, le jeu de données GS-LF, les métriques (AUROC et F1 macro) et des résultats de référence.
La validation prospective de POM montre qu'un modèle appris peut approcher la performance d'un expert humain, ce qui justifie l'intérêt applicatif du sujet.
Enfin, OpenPOM fournit les données nettoyées et le découpage, que nous réutilisons~: notre fichier de données est identique au fichier GS-LF d'OpenPOM.

\subsubsection{Limites des propositions}

La propagation de messages agrège à chaque couche l'information des voisins immédiats d'un atome~: avec $K$ couches, un atome ne <<~voit~>> que les atomes situés à au plus $K$ liaisons.
Les interactions entre atomes éloignés sont donc mal captées.
Surtout, ces modèles n'utilisent que le graphe 2D~: ils ignorent la géométrie de la molécule et, en particulier, sa stéréochimie.
Deux énantiomères, molécules images l'une de l'autre dans un miroir, ont le même graphe et reçoivent nécessairement la même prédiction, alors que leurs odeurs peuvent différer.
Enfin, ces modèles sont entraînés à partir de zéro sur environ 5\,000 molécules, sans bénéficier de connaissances chimiques plus générales.

%--------------------------------------------------------------------------------

\section{Les transformeurs moléculaires 3D~: Uni-Mol et Uni-Mol2}
\label{sec:EAUniMol}

\subsection{Présentation}

\paragraph{Uni-Mol.}
Uni-Mol \cite{Zhou-2023} est un transformeur qui prend en entrée une molécule sous forme 3D~: la liste de ses atomes et leurs coordonnées dans une conformation donnée.
Chaque atome est un \emph{token}, comme un mot dans une phrase, et un token spécial \texttt{[CLS]} est ajouté pour résumer la molécule entière.
La géométrie est encodée par une \emph{représentation de paires}~: pour chaque couple d'atomes $(i, j)$, la distance euclidienne $d_{ij}$ est projetée sur une base de noyaux gaussiens, dont les paramètres dépendent du type de la paire.
Cette représentation de paires est ajoutée comme biais aux scores d'attention et mise à jour à chaque couche.
L'attention étant globale, chaque atome interagit avec tous les autres dès la première couche, quelle que soit leur distance dans le graphe.
Comme seules les distances interviennent, la représentation est invariante par rotation et par translation de la molécule.

Le modèle (15 couches, dimension 512, 64 têtes d'attention, environ 47 millions de paramètres) est pré-entraîné de manière auto-supervisée sur 19 millions de molécules et plus de 200 millions de conformations, avec deux tâches~: retrouver des atomes masqués et débruiter des coordonnées perturbées.
Il est ensuite ajusté (\emph{fine-tuning}) sur des tâches de prédiction de propriétés, où il obtient de très bons résultats.
Une bibliothèque, \texttt{unimol\_tools}, fournit les poids pré-entraînés et les outils de préparation des données.

\paragraph{Uni-Mol2.}
Uni-Mol2 \cite{Ji-2024} étudie le passage à l'échelle de cette approche.
Le jeu de pré-entraînement est porté à environ 884 millions de molécules (73 millions de squelettes moléculaires distincts), et des modèles de 84 millions à 1,1 milliard de paramètres sont entraînés.
L'architecture devient un transformeur <<~à deux pistes~>> qui met à jour conjointement une représentation des atomes et une représentation des paires.
Les paires combinent désormais la géométrie (distances encodées par noyaux gaussiens) et la topologie (type de liaison et plus court chemin dans le graphe).
Les auteurs établissent une loi d'échelle~: la perte de validation décroît en loi de puissance avec la taille du modèle, la quantité de données et le calcul, et les performances sur les tâches en aval progressent avec la taille.

\begin{figure}
   \centering
   \placeholder{schema: architecture Uni-Mol -- atomes + [CLS] en tokens, distances en noyaux gaussiens, représentation de paires utilisée comme biais d'attention, 15 couches, sortie [CLS] de dimension 512}
   \caption{Architecture du transformeur Uni-Mol (d'après \cite{Zhou-2023})}
   \label{fig:UniMol}
\end{figure}

\subsection{Analyse}

\subsubsection{Intérêts des propositions}

Uni-Mol répond directement aux deux limites des GNN identifiées précédemment.
L'attention globale supprime la myopie de la propagation de messages, et la représentation de paires intègre la géométrie, qui conditionne l'interaction d'une molécule odorante avec les récepteurs olfactifs.
Le pré-entraînement sur des centaines de millions de conformations apporte une connaissance chimique générale, précieuse lorsque la tâche cible ne compte que 5\,000 molécules~: c'est le principe de l'apprentissage par transfert.
Enfin, les poids et le code sont publics, et la bibliothèque \texttt{unimol\_tools} rend la version~1 directement utilisable.
Uni-Mol2 est la cible désignée par le sujet~; sa représentation de paires, qui combine géométrie et topologie, devrait mieux exploiter la structure moléculaire.

\subsubsection{Limites des propositions}

\paragraph{Dépendance aux conformères.}
Le modèle requiert des coordonnées 3D, qu'il faut générer (ici avec RDKit).
Une molécule flexible possède de nombreuses conformations~; le choix de celle qui est présentée au modèle est arbitraire et introduit une variabilité.
La génération peut aussi échouer~: c'est le cas de 6 molécules de GS-LF.

\paragraph{Invariance par réflexion.}
Une matrice de distances est inchangée par une symétrie miroir.
Deux énantiomères ont donc exactement les mêmes distances interatomiques~: Uni-Mol ne peut pas les distinguer.
Il ne distingue que les stéréoisomères dont les distances diffèrent, comme les diastéréoisomères.
Uni-Mol2 utilise les mêmes distances, complétées par des informations topologiques~: il semble partager cette limite, ce qui reste à vérifier expérimentalement.

\paragraph{Écart de domaine et coût.}
Les molécules du pré-entraînement proviennent de bases de chimie médicinale et commerciale (ZINC, ChEMBL), dont la distribution diffère de celle des petites molécules odorantes.
Enfin, ces modèles sont nettement plus coûteux que les GNN~: l'attention est quadratique en nombre d'atomes, et Uni-Mol2 compte au moins 84 millions de paramètres.

%--------------------------------------------------------------------------------

\section{Uni-Mol2 appliqué à l'olfaction}
\label{sec:EAHan}

\subsection{Présentation}

Han et al. \cite{Han-2026} étudient si un modèle de fondation moléculaire généraliste, ajusté une seule fois sur une tâche olfactive, apprend une représentation transférable à d'autres tâches olfactives.
Ils ajustent Uni-Mol2 (version 84M, la version 164M n'apportant pas de gain à budget égal) sur GS-LF avec une perte focale, et le comparent à POM et OpenPOM.

\paragraph{Protocole d'évaluation.}
Le protocole est explicitement aligné sur celui de POM, afin de permettre une comparaison équitable (figure~\ref{fig:ProtocoleHan})~:
\begin{enumerate}
   \item environ 20~\% de GS-LF est tenu à l'écart comme ensemble de test, par stratification itérative multi-label~;
   \item une recherche d'hyperparamètres de 50 essais est menée en validation croisée à cinq plis sur le reste~; les 10 meilleures configurations selon l'AUROC de validation sont retenues, soit $10 \times 5 = 50$ modèles~;
   \item un seuil de décision est choisi pour chaque descripteur, en maximisant son F1 sur les prédictions \emph{hors pli} (\emph{out-of-fold})~: chaque molécule d'entraînement est prédite par le modèle qui ne l'a pas vue~;
   \item sur le test, les probabilités des 50 modèles sont moyennées, puis binarisées avec ces seuils~;
   \item les métriques sont des moyennes macro (AUROC, AUPRC, F1, précision, rappel), accompagnées d'intervalles de confiance à 95~\% calculés par 1\,000 rééchantillonnages \emph{bootstrap} du test.
\end{enumerate}

\paragraph{Résultats.}
Sur le test de GS-LF, Uni-Mol2 ajusté dépasse les deux modèles de référence sur les cinq métriques (tableau~\ref{tab:ResultatsLitterature}), avec des écarts modestes et des intervalles de confiance qui se chevauchent.
Les auteurs évaluent ensuite le modèle, sans nouvel entraînement, sur d'autres tâches~: des paires d'énantiomères, un autre jeu de descripteurs, la distinction entre molécules odorantes et inodores, et la discriminabilité de mélanges d'odeurs.
Sur les énantiomères, OpenPOM donne des prédictions strictement identiques aux deux molécules de chaque paire, tandis qu'Uni-Mol2 les différencie~; mais il n'attribue la bonne odeur au bon énantiomère que dans 19 cas sur 34, à peine mieux que le hasard.
Enfin, l'analyse des erreurs montre que les descripteurs confondus par le modèle sont ceux qui apparaissent souvent ensemble, ce qui conduit les auteurs à suggérer de modéliser conjointement molécules et descripteurs.

\begin{table}
   \centering
   \begin{adjustbox}{max width=\linewidth}
   \begin{tabular}{|l||c|c|c|c|c|}
      \hline
      Modèle & AUROC & AUPRC & F1 & Précision & Rappel \\
      \hline
      \hline
      POM (chiffres publiés) & 0,894 & -- & 0,360 & 0,379 & -- \\
      \hline
      OpenPOM & 0,888 & 0,375 & 0,375 & 0,406 & 0,405 \\
      \hline
      Uni-Mol2 ajusté & \textbf{0,899} & \textbf{0,408} & \textbf{0,399} & \textbf{0,429} & \textbf{0,429} \\
      \hline
   \end{tabular}
   \end{adjustbox}
   \caption{Performances sur le test de GS-LF, moyennes macro sur 138 descripteurs (source~: \cite{Han-2026})}
   \label{tab:ResultatsLitterature}
\end{table}

\begin{figure}
   \centering
   \placeholder{schema: protocole d'évaluation de Han et al. -- GS-LF, découpage 80/20 stratifié, validation croisée 5 plis, prédictions hors pli, seuils par descripteur, moyenne des modèles sur le test, bootstrap}
   \caption{Protocole d'évaluation de \cite{Han-2026}, repris de POM}
   \label{fig:ProtocoleHan}
\end{figure}

\subsection{Analyse}

\subsubsection{Intérêts de la proposition}

Ce travail est le plus proche du nôtre~: il utilise le modèle visé par le sujet, sur le même jeu de données, et le compare aux mêmes références.
Il apporte un résultat de référence pour un transformeur 3D sur GS-LF et, surtout, un protocole d'évaluation précis.
Il montre aussi que les performances absolues restent modestes~: un F1 macro autour de 0,40 constitue l'état de l'art, ce qui fixe des attentes réalistes pour notre propre travail.
Enfin, son expérience sur les énantiomères et son analyse des confusions entre descripteurs nourrissent directement nos perspectives.

\subsubsection{Limites de la proposition}

Plusieurs détails nécessaires à la reproduction ne sont pas donnés~: taux d'apprentissage, nombre d'époques, couches dégelées, taille des lots, architecture de la tête de classification et paramètres de la perte focale.
Le code n'était pas accessible au moment de notre lecture.
Le gain sur GS-LF est faible au regard des intervalles de confiance, et l'ensemble de 50 modèles rend la comparaison coûteuse à reproduire.
Enfin, GS-LF ne contenant aucune information stéréochimique, le modèle n'a jamais pu apprendre quel énantiomère sent quoi~: le résultat proche du hasard sur les énantiomères était prévisible.

%--------------------------------------------------------------------------------

\section{Récapitulatif}
\label{sec:Recapitulatif}

Le tableau~\ref{tab:Comparatif} résume les avantages et inconvénients de chaque famille de travaux, et le tableau~\ref{tab:ComparatifAlternatif} les compare selon les critères qui comptent pour notre problème.

\begin{table}
   \centering
   \begin{adjustbox}{max width=\linewidth}
   \begin{tabular}{|l||p{0.36\linewidth}|p{0.36\linewidth}|}
      \hline
      Proposition & Avantages & Inconvénients \\
      \hline
      \hline
      POM / OpenPOM & Référence de la tâche, validation prospective, code et données ouverts (OpenPOM) & Myopie de la propagation de messages, pas de géométrie ni de stéréochimie, pas de pré-entraînement \\
      \hline
      Uni-Mol & Attention globale, géométrie 3D, pré-entraînement massif, outils disponibles & Dépend des conformères générés, invariant par réflexion, coût de calcul \\
      \hline
      Uni-Mol2 & Paires combinant géométrie et topologie, modèles plus grands, loi d'échelle & Coût encore supérieur, mêmes limites géométriques, intégration plus lourde \\
      \hline
      Han et al. & Même tâche et mêmes données, protocole d'évaluation de référence & Détails d'entraînement manquants, gains modestes, ensemble coûteux \\
      \hline
   \end{tabular}
   \end{adjustbox}
   \caption{Tableau comparatif des propositions étudiées}
   \label{tab:Comparatif}
\end{table}

\begin{table*}
   \centering
   \begin{adjustbox}{max width=\linewidth}
   \begin{tabular}{|l||c|c|c|c|}
      \hline
      Critère & POM / OpenPOM & Uni-Mol & Uni-Mol2 & Han et al. \\
      \hline
      \hline
      Interactions entre atomes éloignés & & $\surd$ & $\surd$ & $\surd$ \\
      \hline
      Géométrie 3D & & $\surd$ & $\surd$ & $\surd$ \\
      \hline
      Distinction des énantiomères & & & & $\approx$ \\
      \hline
      Pré-entraînement à grande échelle & & $\surd$ & $\surd$ & $\surd$ \\
      \hline
      Évalué sur GS-LF & $\surd$ & & & $\surd$ \\
      \hline
      Protocole d'évaluation reproductible & $\surd$ & & & $\approx$ \\
      \hline
      Code et poids disponibles & $\surd$ & $\surd$ & $\surd$ & \\
      \hline
      Coût de calcul modéré & $\surd$ & $\approx$ & & \\
      \hline
   \end{tabular}
   \end{adjustbox}
   \caption{Tableau comparatif des avantages des propositions étudiées ($\surd$~: oui~; $\approx$~: partiellement)}
   \label{tab:ComparatifAlternatif}
\end{table*}

%--------------------------------------------------------------------------------

\section{Conclusion}
\label{sec:ConclusionEtatArt}

Les GNN de référence définissent le cadre et les performances à atteindre sur GS-LF, mais n'exploitent ni la géométrie ni un pré-entraînement à grande échelle.
Uni-Mol et Uni-Mol2 apportent ces deux éléments, et le travail de Han et al. confirme qu'Uni-Mol2 ajusté atteint, voire dépasse légèrement, l'état de l'art sur cette tâche.

Trois directions se dégagent pour notre travail~:
\begin{enumerate}
   \item adapter un transformeur Uni-Mol à la classification multi-label des 138 descripteurs, avec un décodeur MLP, en commençant par la version~1, directement disponible, avant de passer à Uni-Mol2~;
   \item traiter explicitement le déséquilibre des descripteurs et la faible quantité de données~;
   \item reproduire le protocole d'évaluation de \cite{Han-2026}, condition nécessaire pour comparer nos résultats à ceux de la littérature.
\end{enumerate}
Ces directions font l'objet du chapitre suivant.

\chapter{Proposition}
\label{chap:Propositions}

Partant de l'hypothèse qu'une représentation 3D pré-entraînée capture mieux les déterminants structurels de l'odeur qu'un GNN 2D appris à partir de zéro, nous proposons d'adapter le transformeur Uni-Mol à la classification multi-label des 138 descripteurs de GS-LF.
Cette proposition se décompose en trois sous-problèmes complémentaires~: la représentation des molécules et l'encodage par le transformeur, l'apprentissage du décodeur multi-label malgré le déséquilibre, et l'évaluation selon un protocole comparable à la littérature.
La section~\ref{sec:PropositionUniMol} présente successivement les idées qui la fondent, les outils nécessaires, sa formalisation et son analyse.

%--------------------------------------------------------------------------------

\section{Adaptation d'Uni-Mol à la classification multi-label des odeurs}
\label{sec:PropositionUniMol}

\subsection{Idées préliminaires}
\label{sec:IdeesPreliminaires}

\paragraph{Transférer plutôt qu'apprendre.}
Avec environ 5\,000 molécules, entraîner un transformeur de 47 millions de paramètres à partir de zéro conduirait à un surapprentissage massif.
Nous réutilisons donc Uni-Mol pré-entraîné comme \emph{encodeur}~: il transforme une molécule en un vecteur de dimension 512 (la sortie du token \texttt{[CLS]}), sur lequel un petit décodeur MLP prédit les 138 descripteurs.
Deux modes d'apprentissage sont envisagés~: encodeur entièrement gelé, où seul le décodeur apprend, puis dégel des dernières couches de l'encodeur, ajustées avec un taux d'apprentissage plus faible que celui du décodeur.

\paragraph{Utiliser la géométrie, et sa variabilité.}
Une molécule n'a pas une forme unique mais un ensemble de conformations.
Plutôt que d'en choisir une, nous en générons plusieurs par molécule et en tirons une au hasard à chaque passage pendant l'entraînement.
Ce tirage, complété par un léger bruit sur les coordonnées, constitue une augmentation de données naturelle~: le modèle apprend une représentation robuste aux variations de conformation, ce qui est précieux lorsque les données sont rares.

\paragraph{Compenser le déséquilibre.}
La plupart des descripteurs sont absents de la plupart des molécules~: une perte d'entropie croisée binaire standard est dominée par ces négatifs faciles, et le modèle apprend à prédire <<~absent~>> partout.
Nous comparons deux remèdes~: pondérer les positifs de chaque descripteur par l'inverse de leur fréquence, et utiliser une perte asymétrique \cite{Ridnik-2021}, qui atténue la contribution des négatifs faciles.

\paragraph{Évaluer comme la littérature.}
Les métriques n'ont de sens que relativement à un protocole.
Nous adoptons celui de \cite{Han-2026}, décrit à la section~\ref{sec:EAHan}, pour que nos résultats soient directement comparables à ceux de POM, OpenPOM et Uni-Mol2.

\subsection{Outils nécessaires}
\label{sec:OutilsNecessaires}

La proposition s'appuie sur les briques suivantes, détaillées au chapitre~\ref{chap:Experimentations}~:
\begin{itemize}
   \item \textbf{RDKit}\footnote{\url{https://www.rdkit.org}}, pour lire les codes SMILES, ajouter les hydrogènes et générer des conformations 3D~;
   \item \textbf{les poids pré-entraînés d'Uni-Mol} et la bibliothèque \texttt{unimol\_tools}, qui fournit l'architecture et le dictionnaire des atomes~;
   \item \textbf{PyTorch} et \textbf{PyTorch Lightning}, pour le décodeur, l'apprentissage et le calcul sur GPU~;
   \item \textbf{scikit-learn} et \textbf{scikit-multilearn}, pour les métriques et la stratification itérative multi-label \cite{Sechidis-2011, Szymanski-2017}.
\end{itemize}

\subsection{Formalisation}
\label{sec:Formalisation}

\subsubsection{Définition du problème}

Soit $\mathcal{D} = \{(m_i, \mathbf{y}_i)\}_{i=1}^{N}$ le jeu de données, où $m_i$ est une molécule décrite par son code SMILES et $\mathbf{y}_i \in \{0, 1\}^{L}$ le vecteur de ses descripteurs, avec $N = 4\,983$ et $L = 138$~: $y_{ij} = 1$ si le descripteur $j$ est attribué à la molécule $i$.
On cherche une fonction $f_\theta$ qui associe à une molécule un vecteur de probabilités $\hat{\mathbf{p}} = f_\theta(m) \in [0, 1]^{L}$, puis des seuils $\boldsymbol{\tau} \in [0, 1]^{L}$ tels que la prédiction $\hat{y}_j = \mathbb{1}[\hat{p}_j \geq \tau_j]$ soit la plus proche possible de $y_j$ pour chaque descripteur.

\subsubsection{Représentation 3D}

Pour chaque molécule $m$, on génère un ensemble de conformations $\mathcal{C}(m) = \{\mathbf{X}^{(1)}, \ldots, \mathbf{X}^{(K_m)}\}$, où $\mathbf{X}^{(k)} \in \mathbb{R}^{n \times 3}$ contient les coordonnées des $n$ atomes, hydrogènes compris.
Une conformation est encodée par~:
\begin{itemize}
   \item la suite des tokens $(t_0, t_1, \ldots, t_n, t_{n+1})$, où $t_0 = \texttt{[CLS]}$, $t_{n+1} = \texttt{[SEP]}$ et $t_1, \ldots, t_n$ sont les indices des atomes dans le dictionnaire d'Uni-Mol, de taille $V = 31$~;
   \item les coordonnées centrées $\mathbf{x}_k - \bar{\mathbf{x}}$, les tokens spéciaux étant placés à l'origine~;
   \item la matrice des distances $d_{kl} = \lVert \mathbf{x}_k - \mathbf{x}_l \rVert_2$~;
   \item le type de chaque paire, $e_{kl} = t_k \times V + t_l$, qui permet au modèle d'adapter l'encodage des distances aux éléments chimiques concernés.
\end{itemize}

\subsubsection{Encodeur}

Uni-Mol encode chaque distance par $K = 128$ noyaux gaussiens dont l'échelle et le décalage dépendent du type de paire~:
\begin{equation}
   \label{eq:Gaussiennes}
   \boldsymbol{\psi}_{kl} = \left( \mathcal{G}\left(a_{e_{kl}} \, d_{kl} + b_{e_{kl}} ;\ \mu_r, \sigma_r\right) \right)_{r = 1, \ldots, K},
   \qquad
   \mathcal{G}(x ; \mu, \sigma) = \frac{1}{\sigma \sqrt{2 \pi}} \, e^{-\frac{(x - \mu)^2}{2 \sigma^2}}
\end{equation}
Une projection linéaire de $\boldsymbol{\psi}_{kl}$ initialise la représentation de paires $\mathbf{q}^{0}_{kl}$, qui biaise l'attention de chaque couche $\ell$~:
\begin{equation}
   \label{eq:Attention}
   \mathrm{Attention}(\mathbf{Q}, \mathbf{K}, \mathbf{V}) = \mathrm{softmax}\left( \frac{\mathbf{Q} \mathbf{K}^{\top}}{\sqrt{d}} + \mathbf{q}^{\ell} \right) \mathbf{V},
   \qquad
   \mathbf{q}^{\ell + 1}_{kl} = \mathbf{q}^{\ell}_{kl} + \frac{\mathbf{Q}_k \mathbf{K}_l^{\top}}{\sqrt{d}}
\end{equation}
Après les 15 couches, la sortie du token \texttt{[CLS]} fournit la représentation de la molécule $\mathbf{h} = g_\phi(\mathbf{X}) \in \mathbb{R}^{512}$.
Les paramètres $\phi$ de l'encodeur sont gelés, à l'exception éventuelle des $u$ dernières couches.
Comme la représentation ne dépend que des distances, elle est invariante par toute isométrie de l'espace, rotations, translations \emph{et réflexions}.

\subsubsection{Décodeur multi-label}

Le décodeur est un perceptron multicouche à une couche cachée de dimension 256~:
\begin{equation}
   \label{eq:Decodeur}
   \mathbf{z} = \mathbf{W}_2 \, \mathrm{Dropout}\big(\mathrm{GELU}(\mathrm{LN}(\mathbf{W}_1 \mathbf{h} + \mathbf{b}_1))\big) + \mathbf{b}_2,
   \qquad
   \hat{\mathbf{p}} = \sigma(\mathbf{z})
\end{equation}
avec $\mathbf{W}_1 \in \mathbb{R}^{256 \times 512}$, $\mathbf{W}_2 \in \mathbb{R}^{138 \times 256}$, $\mathrm{LN}$ une normalisation de couche et $\sigma$ la fonction sigmoïde appliquée à chaque logit $z_j$ indépendamment.
Chaque ligne $\mathbf{w}_j$ de $\mathbf{W}_2$ est le vecteur propre au descripteur $j$~: le logit vaut $z_j = \mathbf{w}_j \cdot \mathbf{h}' + b_j$, où $\mathbf{h}'$ est la sortie de la couche cachée.
Les 138 descripteurs sont ainsi traités comme des tâches binaires indépendantes partageant la même représentation.

\subsubsection{Fonction de perte}

\paragraph{Entropie croisée pondérée.}
Pour un lot $\mathcal{B}$, la perte de référence est l'entropie croisée binaire dont les positifs du descripteur $j$ sont pondérés par $w_j = \min\left(\frac{N - N_j^{+}}{N_j^{+}},\ w_{\max}\right)$, où $N_j^{+}$ est le nombre de positifs du descripteur $j$ dans l'ensemble d'entraînement et $w_{\max} = 20$ un plafond qui évite des gradients excessifs pour les descripteurs les plus rares~:
\begin{equation}
   \label{eq:BCE}
   \mathcal{L}_{\mathrm{BCE}} = -\frac{1}{|\mathcal{B}|\,L} \sum_{i \in \mathcal{B}} \sum_{j=1}^{L} \Big[ w_j \, y_{ij} \log \hat{p}_{ij} + (1 - y_{ij}) \log (1 - \hat{p}_{ij}) \Big]
\end{equation}

\paragraph{Perte asymétrique.}
La perte asymétrique (ASL) \cite{Ridnik-2021} traite différemment positifs et négatifs.
Elle décale d'abord la probabilité des négatifs d'une marge $c$, $\hat{p}^{\,c}_{ij} = \max(\hat{p}_{ij} - c, 0)$, ce qui annule la perte des négatifs très faciles, puis les pondère par un facteur focal d'exposant $\gamma_-$~:
\begin{equation}
   \label{eq:ASL}
   \mathcal{L}_{\mathrm{ASL}} = -\frac{1}{|\mathcal{B}|} \sum_{i \in \mathcal{B}} \sum_{j=1}^{L} \Big[ y_{ij} \, (1 - \hat{p}_{ij})^{\gamma_+} \log \hat{p}_{ij} + (1 - y_{ij}) \, \big(\hat{p}^{\,c}_{ij}\big)^{\gamma_-} \log \big(1 - \hat{p}^{\,c}_{ij}\big) \Big]
\end{equation}
Nous utilisons $\gamma_+ = 0$, $\gamma_- = 2$ et $c = 0{,}05$~: les positifs, rares, conservent tout leur poids, tandis que les négatifs déjà bien classés contribuent peu.
Les facteurs focaux sont traités comme des constantes lors du calcul du gradient, comme dans l'implémentation de référence.

\subsubsection{Augmentation géométrique}

À chaque passage d'une molécule $m$ pendant l'entraînement, on tire une conformation $\mathbf{X}$ uniformément parmi les conformations relaxées de $\mathcal{C}(m)$, puis on perturbe ses coordonnées~:
\begin{equation}
   \label{eq:Jitter}
   \tilde{\mathbf{x}}_k = \mathbf{x}_k + \boldsymbol{\varepsilon}_k,
   \qquad
   \boldsymbol{\varepsilon}_k \sim \mathcal{N}(\mathbf{0}, s^2 \mathbf{I}_3),
   \qquad
   s \sim \mathcal{U}([s_{\min}, s_{\max}])
\end{equation}
avec $s_{\min} = 0{,}01$~\AA{} et $s_{\max} = 0{,}04$~\AA, des amplitudes faibles devant la longueur d'une liaison covalente (de l'ordre de 1~\AA).
En validation et en test, la conformation est fixe et aucun bruit n'est ajouté.
Une permutation aléatoire de l'ordre des atomes n'est pas utile~: l'attention sans encodage de position rend déjà Uni-Mol invariant à cet ordre.

\subsubsection{Protocole d'évaluation}

\paragraph{Découpage.}
GS-LF est découpé en un ensemble d'entraînement $\mathcal{D}_{\mathrm{train}}$ (80~\%) et un ensemble de test $\mathcal{D}_{\mathrm{test}}$ (20~\%) par stratification itérative d'ordre~2 \cite{Sechidis-2011, Szymanski-2017}, qui préserve la fréquence de chaque descripteur et de chaque paire de descripteurs.
$\mathcal{D}_{\mathrm{train}}$ est à son tour découpé en cinq plis $F_1, \ldots, F_5$ de la même manière.

\paragraph{Seuils hors pli.}
Pour chaque pli $k$, un modèle $f_k$ est entraîné sur $\mathcal{D}_{\mathrm{train}} \setminus F_k$, sa meilleure époque étant choisie selon l'AUROC macro sur $F_k$.
Chaque molécule d'entraînement reçoit ainsi une prédiction hors pli $\hat{\mathbf{p}}^{\mathrm{oof}}_i = f_{k(i)}(m_i)$, faite par le seul modèle qui ne l'a pas vue.
Le seuil de chaque descripteur maximise son F1 sur ces prédictions~:
\begin{equation}
   \label{eq:Seuils}
   \tau_j = \underset{\tau}{\arg\max} \ \mathrm{F1}\Big( \big(y_{ij}\big)_{i \in \mathcal{D}_{\mathrm{train}}},\ \big(\mathbb{1}[\hat{p}^{\mathrm{oof}}_{ij} \geq \tau]\big)_{i \in \mathcal{D}_{\mathrm{train}}} \Big)
\end{equation}

\paragraph{Ensemble et métriques.}
Sur le test, les probabilités des cinq modèles sont moyennées, $\hat{\mathbf{p}}_i = \frac{1}{5} \sum_{k=1}^{5} f_k(m_i)$, puis binarisées avec les seuils $\tau_j$.
Les métriques sont des moyennes macro, c'est-à-dire des moyennes non pondérées sur les descripteurs, ce qui donne le même poids aux descripteurs rares qu'aux fréquents~:
\begin{equation}
   \label{eq:F1Macro}
   \mathrm{F1}_{\mathrm{macro}} = \frac{1}{L} \sum_{j=1}^{L} \frac{2 \, P_j \, R_j}{P_j + R_j}
\end{equation}
où $P_j$ et $R_j$ sont la précision et le rappel du descripteur $j$.
L'AUROC et l'AUPRC macro, indépendantes des seuils, sont calculées de la même façon sur les descripteurs présents dans le test.
Enfin, l'incertitude est estimée en recalculant les métriques sur 1\,000 rééchantillonnages avec remise du test~; l'intervalle de confiance à 95~\% est donné par les percentiles 2,5 et 97,5.

\begin{figure}
   \centering
   \placeholder{schema: chaîne complète -- SMILES, RDKit (conformères), encodage Uni-Mol (tokens, distances, types de paires), augmentation, encodeur Uni-Mol, [CLS] 512, décodeur MLP, 138 logits, sigmoïde, seuils par descripteur}
   \caption{Vue d'ensemble de la chaîne proposée}
   \label{fig:Pipeline}
\end{figure}

\subsection{Analyse}
\label{sec:AnalyseProposition}

\paragraph{Ce que la proposition résout.}
Elle répond aux limites des GNN identifiées au chapitre~\ref{chap:EtatArt}~: l'attention globale capte les interactions entre atomes éloignés, la représentation de paires intègre la géométrie, et le pré-entraînement compense en partie la rareté des données.
La perte asymétrique et les seuils par descripteur traitent le déséquilibre à deux niveaux distincts, l'apprentissage et la décision.
Le protocole d'évaluation rend la comparaison avec la littérature possible.

\paragraph{Ce qu'elle ne résout pas.}
L'invariance par réflexion (équation~\ref{eq:Attention} et suivantes) empêche toute distinction entre énantiomères, et GS-LF, dépourvu de stéréochimie, ne permettrait de toute façon pas de l'apprendre.
Les descripteurs restent traités comme indépendants, alors que leurs co-occurrences sont fortes et que les confusions du modèle en dépendent \cite{Han-2026}.
Enfin, la variabilité des annotations humaines borne la performance atteignable, quel que soit le modèle.

\paragraph{Coût de calcul.}
L'attention est quadratique en nombre d'atomes~: avec les hydrogènes, les molécules de GS-LF comptent jusqu'à 184 atomes.
Le dégel de couches impose de rétropropager à travers l'encodeur, ce qui contraint la taille des lots~; nous utilisons des lots de 8 avec accumulation de gradient sur 8 lots.
Le protocole complet multiplie le coût d'un entraînement par cinq, contre cinquante pour celui de Han et al.

%--------------------------------------------------------------------------------

\section{Conclusion}
\label{sec:ConclusionProposition}

Notre proposition combine un encodeur 3D pré-entraîné, une augmentation par conformations, un décodeur MLP entraîné avec une perte asymétrique et un dégel partiel de l'encodeur, ainsi qu'un protocole d'évaluation aligné sur la littérature.
Plusieurs de ses composantes sont des choix à valider empiriquement~: le mode d'apprentissage, la fonction de perte, l'intérêt de l'augmentation, la profondeur du décodeur ou la stratégie de découpage.
Le chapitre suivant décrit leur mise en œuvre et les expériences menées pour les départager.

\chapter{Expérimentations et résultats}
\label{chap:Experimentations}

Ce chapitre présente d'abord les outils utilisés et les développements réalisés pour mettre en œuvre la proposition du chapitre~\ref{chap:Propositions}, ainsi que les contraintes rencontrées.
Il décrit ensuite la mise en place des expériences, leurs résultats et leur analyse.

%--------------------------------------------------------------------------------

\section{Outils et frameworks utilisés}
\label{sec:Outils}

Le tableau~\ref{tab:Outils} récapitule les outils utilisés.
Le code est écrit en Python et organisé selon un modèle de projet associant PyTorch Lightning, qui structure l'apprentissage (boucles d'entraînement et de validation, points de sauvegarde, précision numérique), et Hydra, qui compose la configuration à partir de fichiers YAML.
Chaque expérience est ainsi décrite par un unique fichier de configuration, qui ne contient que ses écarts à la configuration par défaut~; cela rend les expériences reproductibles et faciles à comparer.

\begin{table}
   \centering
   \begin{adjustbox}{max width=\linewidth}
   \begin{tabular}{|l||p{0.62\linewidth}|}
      \hline
      Outil & Usage \\
      \hline
      \hline
      PyTorch, PyTorch Lightning & Modèle, apprentissage, calcul sur GPU \\
      \hline
      Hydra & Configuration composable des expériences \\
      \hline
      RDKit & Lecture des SMILES, génération et relaxation des conformations 3D, squelettes de Bemis-Murcko \\
      \hline
      \texttt{unimol\_tools} & Architecture et poids pré-entraînés d'Uni-Mol, dictionnaire des atomes \\
      \hline
      scikit-learn, TorchMetrics & Métriques (AUROC, AUPRC, F1, précision, rappel) \\
      \hline
      scikit-multilearn, iterstrat & Découpages stratifiés multi-label \\
      \hline
      TensorBoard & Suivi des courbes d'apprentissage \\
      \hline
      uv, ruff, Git & Environnement reproductible, qualité du code, gestion de versions \\
      \hline
   \end{tabular}
   \end{adjustbox}
   \caption{Outils et bibliothèques utilisés}
   \label{tab:Outils}
\end{table}

Les expériences ont été exécutées sur deux types de machines~: un ordinateur Apple à puce M (accélération MPS) et un poste Windows équipé d'un GPU Intel Arc (accélération XPU)~; \acompleter{modèles exacts des machines et mémoire}.

%--------------------------------------------------------------------------------

\section{Développement et algorithmes implémentés}
\label{sec:Developpement}

\subsection{Architecture logicielle}

Le code est organisé en quatre composants, représentés sur la figure~\ref{fig:Architecture}~:
\begin{itemize}
   \item un \textbf{script hors ligne} qui génère les conformations de chaque molécule et les sauvegarde (algorithme~\ref{alg:Conformeres})~;
   \item un \textbf{module de données} qui découpe le jeu de données, applique l'augmentation, encode chaque conformation au format Uni-Mol (algorithme~\ref{alg:Exemple}) et assemble les lots~;
   \item un \textbf{module d'apprentissage} qui associe l'encodeur Uni-Mol, le décodeur, la fonction de perte, l'optimiseur et les métriques (algorithme~\ref{alg:Entrainement})~;
   \item un \textbf{script de validation croisée} qui enchaîne les plis et calcule les métriques du protocole d'évaluation (algorithme~\ref{alg:Evaluation}).
\end{itemize}

\begin{figure}
   \centering
   \placeholder{schema: architecture logicielle -- script de génération des conformères, module de données (découpage, augmentation, encodage, assemblage des lots), module d'apprentissage (encodeur, décodeur, perte, métriques), script de validation croisée, configurations Hydra, TensorBoard}
   \caption{Architecture logicielle de la chaîne d'apprentissage}
   \label{fig:Architecture}
\end{figure}

\subsection{Génération des conformations}

Les conformations sont générées une fois pour toutes, avant l'entraînement (algorithme~\ref{alg:Conformeres}).
L'algorithme ETKDGv3 de RDKit produit jusqu'à 15 géométries compatibles avec la topologie, en écartant celles trop proches d'une géométrie déjà retenue (écart quadratique moyen inférieur à 0,25~\AA).
Ces géométries sont ensuite relaxées par le champ de force MMFF94, qui les rapproche d'un minimum d'énergie.
Un tiers des géométries brutes est conservé en plus des géométries relaxées.

Sur les 4\,983 molécules, 4\,977 ont pu être traitées~; la génération a échoué pour 6 molécules, qui sont exclues de l'apprentissage.
On obtient 71\,435 conformations, entre 2 et 20 par molécule (14,4 en moyenne), pour des molécules de 1 à 184 atomes, hydrogènes compris.

\begin{algorithm}
   \caption{Génération hors ligne des conformations}
   \label{alg:Conformeres}
   \begin{algorithmic}[1]
      \REQUIRE jeu de données $\{(s_i, \mathbf{y}_i)\}$ (SMILES et descripteurs), $K_{\max} = 15$, seuil de diversité $\delta = 0{,}25$~\AA
      \ENSURE pour chaque molécule, ses atomes et une liste de conformations
      \FORALL{molécule $(s_i, \mathbf{y}_i)$}
         \STATE $m \leftarrow$ graphe moléculaire de $s_i$, hydrogènes explicites ajoutés
         \STATE $\mathcal{C} \leftarrow$ ETKDGv3$(m, K_{\max}, \delta)$ \COMMENT{géométries diverses}
         \IF{$\mathcal{C} = \emptyset$}
            \STATE $\mathcal{C} \leftarrow$ ETKDGv3 à coordonnées initiales aléatoires, sans élagage
         \ENDIF
         \IF{$\mathcal{C} = \emptyset$}
            \STATE enregistrer l'échec et passer à la molécule suivante
         \ENDIF
         \STATE $\mathcal{C}_{\mathrm{brut}} \leftarrow$ premier tiers de $\mathcal{C}$
         \STATE $\mathcal{C}_{\mathrm{relax}} \leftarrow$ MMFF94$(\mathcal{C})$ \COMMENT{relaxation, au plus 200 itérations}
         \STATE sauvegarder (atomes de $m$, $\mathcal{C}_{\mathrm{brut}} \cup \mathcal{C}_{\mathrm{relax}}$, $\mathbf{y}_i$)
      \ENDFOR
   \end{algorithmic}
\end{algorithm}

\subsection{Préparation des exemples}

À chaque appel, le module de données construit les tenseurs attendus par Uni-Mol à partir d'une conformation (algorithme~\ref{alg:Exemple}).
Les molécules d'un lot n'ayant pas le même nombre d'atomes, les tenseurs sont complétés jusqu'à la taille de la plus grande molécule du lot par le token de remplissage \texttt{[PAD]}, ignoré par l'attention.

\begin{algorithm}
   \caption{Préparation d'un exemple au format Uni-Mol}
   \label{alg:Exemple}
   \begin{algorithmic}[1]
      \REQUIRE molécule (atomes $a_1, \ldots, a_n$, conformations $\mathcal{C}$, descripteurs $\mathbf{y}$), dictionnaire de taille $V = 31$, mode (entraînement ou évaluation)
      \ENSURE tokens, coordonnées, distances, types de paires, descripteurs
      \IF{entraînement}
         \STATE $\mathbf{X} \leftarrow$ conformation tirée uniformément parmi les conformations relaxées de $\mathcal{C}$
         \STATE $s \sim \mathcal{U}([0{,}01 ; 0{,}04])$~; $\mathbf{X} \leftarrow \mathbf{X} + \boldsymbol{\varepsilon}$ avec $\boldsymbol{\varepsilon} \sim \mathcal{N}(\mathbf{0}, s^2 \mathbf{I})$
      \ELSE
         \STATE $\mathbf{X} \leftarrow$ dernière conformation de $\mathcal{C}$ \COMMENT{déterministe}
      \ENDIF
      \STATE $\mathbf{X} \leftarrow \mathbf{X} - \bar{\mathbf{X}}$ \COMMENT{centrage}
      \STATE $\mathbf{t} \leftarrow (\texttt{[CLS]}, \mathrm{idx}(a_1), \ldots, \mathrm{idx}(a_n), \texttt{[SEP]})$
      \STATE ajouter l'origine comme coordonnées de \texttt{[CLS]} et \texttt{[SEP]}
      \STATE $D_{kl} \leftarrow \lVert \mathbf{x}_k - \mathbf{x}_l \rVert_2$ pour tout couple $(k, l)$
      \STATE $E_{kl} \leftarrow t_k \times V + t_l$ pour tout couple $(k, l)$
      \RETURN $(\mathbf{t}, \mathbf{X}, D, E, \mathbf{y})$
   \end{algorithmic}
\end{algorithm}

\subsection{Apprentissage}

L'algorithme~\ref{alg:Entrainement} décrit une époque d'apprentissage.
Les paramètres sont répartis en deux groupes~: le décodeur, appris avec le taux $\eta$, et les couches dégelées de l'encodeur, apprises avec un taux plus faible $\eta_{\mathrm{enc}}$ afin de ne pas détruire les connaissances acquises au pré-entraînement.
Les couches gelées sont maintenues en mode évaluation, ce qui désactive leur \emph{dropout}.
Les gradients de 8 lots de 8 molécules sont accumulés avant chaque mise à jour, ce qui équivaut à des lots de 64 molécules sans saturer la mémoire du GPU.

\begin{algorithm}
   \caption{Une époque d'apprentissage}
   \label{alg:Entrainement}
   \begin{algorithmic}[1]
      \REQUIRE encodeur $g_\phi$ dont les $u$ dernières couches sont dégelées, décodeur $h_\theta$, taux $\eta$ et $\eta_{\mathrm{enc}}$, accumulation $A = 8$
      \ENSURE paramètres mis à jour, métriques de validation
      \STATE couches gelées de $g_\phi$ en mode évaluation~; couches dégelées et $h_\theta$ en mode entraînement
      \FORALL{lot $b$ de l'ensemble d'entraînement (algorithme~\ref{alg:Exemple})}
         \STATE $\mathbf{z} \leftarrow h_\theta(g_\phi(b))$ \COMMENT{logits, $|b| \times 138$}
         \STATE $\ell \leftarrow \mathcal{L}_{\mathrm{ASL}}(\mathbf{z}, \mathbf{y}_b) / A$ \COMMENT{équation~\ref{eq:ASL}}
         \STATE rétropropager $\ell$
         \IF{$A$ lots ont été accumulés}
            \STATE écrêter la norme du gradient à 1
            \STATE mise à jour AdamW~: $\theta$ avec le taux $\eta$, couches dégelées de $\phi$ avec $\eta_{\mathrm{enc}}$
            \STATE remettre les gradients à zéro
         \ENDIF
      \ENDFOR
      \STATE ajuster $\eta$ et $\eta_{\mathrm{enc}}$ selon la décroissance cosinus
      \STATE prédire l'ensemble de validation~; calculer AUROC et F1 macro
      \STATE sauvegarder le modèle si la métrique surveillée s'améliore
   \end{algorithmic}
\end{algorithm}

\subsection{Protocole d'évaluation}

Le protocole de la section~\ref{sec:Formalisation} est implémenté en deux temps.
Un script calcule une fois pour toutes le découpage (1\,007 molécules de test, 3\,976 d'entraînement réparties en cinq plis) et l'enregistre dans un fichier, ce qui garantit que tous les modèles sont évalués sur les mêmes molécules.
Le script de validation croisée applique ensuite l'algorithme~\ref{alg:Evaluation}.

\begin{algorithm}
   \caption{Protocole d'évaluation (d'après \cite{Han-2026})}
   \label{alg:Evaluation}
   \begin{algorithmic}[1]
      \REQUIRE jeu de données GS-LF, nombre de plis $P = 5$, nombre de tirages $B = 1\,000$
      \ENSURE métriques macro sur le test et intervalles de confiance à 95~\%
      \STATE $(\mathcal{D}_{\mathrm{train}}, \mathcal{D}_{\mathrm{test}}) \leftarrow$ stratification itérative d'ordre 2 (80~\% / 20~\%)
      \STATE $(F_1, \ldots, F_P) \leftarrow$ stratification itérative d'ordre 2 de $\mathcal{D}_{\mathrm{train}}$
      \FOR{$k = 1$ \TO $P$}
         \STATE entraîner $f_k$ sur $\mathcal{D}_{\mathrm{train}} \setminus F_k$, en retenant l'époque de meilleure AUROC sur $F_k$
         \STATE $\hat{\mathbf{p}}^{\mathrm{oof}}_i \leftarrow f_k(m_i)$ pour tout $m_i \in F_k$
         \STATE $\hat{\mathbf{p}}^{(k)}_i \leftarrow f_k(m_i)$ pour tout $m_i \in \mathcal{D}_{\mathrm{test}}$
      \ENDFOR
      \FOR{$j = 1$ \TO $138$}
         \STATE $\tau_j \leftarrow$ seuil maximisant le F1 du descripteur $j$ sur $\hat{\mathbf{p}}^{\mathrm{oof}}$ \COMMENT{équation~\ref{eq:Seuils}}
      \ENDFOR
      \STATE $\hat{\mathbf{p}}_i \leftarrow \frac{1}{P} \sum_{k} \hat{\mathbf{p}}^{(k)}_i$ pour tout $m_i \in \mathcal{D}_{\mathrm{test}}$ \COMMENT{ensemble}
      \STATE $M \leftarrow$ métriques macro de $(\hat{\mathbf{p}}, \boldsymbol{\tau})$ sur $\mathcal{D}_{\mathrm{test}}$
      \FOR{$b = 1$ \TO $B$}
         \STATE $M_b \leftarrow$ métriques macro sur un rééchantillonnage avec remise de $\mathcal{D}_{\mathrm{test}}$
      \ENDFOR
      \RETURN $M$ et les percentiles 2,5 et 97,5 des $M_b$
   \end{algorithmic}
\end{algorithm}

\subsection{Contraintes rencontrées}

\subsubsection{Contraintes liées à la proposition}

L'encodage d'entrée doit reproduire \emph{exactement} celui utilisé au pré-entraînement~: le modèle n'a appris à interpréter que celui-là.
Or rien ne signale une erreur d'encodage~: le modèle s'entraîne normalement, mais sur des entrées qu'il interprète mal.
Trois écarts ont ainsi été détectés et corrigés au début de novembre~:
\begin{itemize}
   \item le dictionnaire de la bibliothèque \texttt{unimol\_\allowbreak tools} ne contient pas le token \texttt{[MASK]}, ajouté au pré-entraînement~; sans lui, le vocabulaire compte 30 entrées au lieu de 31, et les types de paires $t_k \times V + t_l$ ne correspondent plus à ceux appris~;
   \item les distances entre le token \texttt{[CLS]} et les atomes n'étaient pas calculées comme au pré-entraînement~;
   \item les couches gelées de l'encodeur restaient en mode entraînement, ce qui laissait leur \emph{dropout} actif et bruitait la représentation.
\end{itemize}
Les expériences antérieures à ces corrections (8 à 15) en ont été affectées.

\subsubsection{Contraintes liées au développement}

\paragraph{Matériel hétérogène.}
L'utilisation d'un GPU Intel Arc sous Windows a nécessité plusieurs contournements dans les pilotes~: transferts mémoire synchrones (des transferts asynchrones corrompaient les étiquettes), désactivation des noyaux vectorisés de l'optimiseur et de l'écrêtage du gradient, et synchronisations explicites entre la propagation et la mise à jour.
Le code isole ces adaptations dans un module dédié, afin de fonctionner indifféremment sur CUDA, MPS, XPU et CPU.

\paragraph{Coût de calcul.}
L'encodage à la volée et l'augmentation empêchent de précalculer les représentations, et le dégel de couches impose de rétropropager à travers l'encodeur.
La durée des expériences (\acompleter{durées mesurées par époque et par expérience}) a limité leur nombre et interdit, à ce stade, une recherche systématique d'hyperparamètres.

%--------------------------------------------------------------------------------

\section{Mise en place des expérimentations}
\label{sec:MiseEnPlace}

\paragraph{Démarche.}
Les expériences ont été conduites de manière incrémentale~: chacune modifie un élément de la précédente pour en mesurer l'effet, et est consignée avec son fichier de configuration, ses courbes et un commentaire.
Les expériences 11 et 12 sont des ablations, qui retirent un élément pour vérifier son utilité.

\paragraph{Évaluation interne.}
Avant l'implémentation du protocole de la section~\ref{sec:Formalisation}, les modèles ont été évalués sur un unique découpage, au niveau des molécules (aucune conformation d'une molécule d'entraînement n'est en validation)~:
\begin{itemize}
   \item expériences 1 à 7~: découpage aléatoire 80~/~10~/~10~;
   \item expériences 8 à 17~: découpage stratifié multi-label 80~/~20, sauf l'expérience 10 qui regroupe les molécules par squelette de Bemis-Murcko, afin qu'aucune famille chimique ne soit à la fois en entraînement et en validation.
\end{itemize}
Les métriques suivantes sont calculées sur la validation à chaque époque~:
\begin{itemize}
   \item l'AUROC macro, indépendante de tout seuil~;
   \item le F1 macro au seuil de 0,5~;
   \item à partir de l'expérience 9, le F1 macro au meilleur seuil \emph{global}, choisi parmi 19 valeurs entre 0,05 et 0,95~;
   \item à partir de l'expérience 14, le F1 micro au meilleur seuil global, et à partir de l'expérience 16, le F1 macro avec un seuil optimal \emph{par descripteur}.
\end{itemize}
Les valeurs rapportées sont les meilleures valeurs atteintes au fil des époques.

\paragraph{Configuration retenue.}
Le tableau~\ref{tab:Hyperparametres} donne la configuration de la meilleure expérience (17) et celle préparée pour le protocole d'évaluation.

\begin{table}
   \centering
   \begin{adjustbox}{max width=\linewidth}
   \begin{tabular}{|l||c|c|}
      \hline
      Paramètre & Expérience 17 & Protocole d'évaluation \\
      \hline
      \hline
      Encodeur & \multicolumn{2}{c|}{Uni-Mol v1, 2 dernières couches dégelées} \\
      \hline
      Décodeur & \multicolumn{2}{c|}{$512 \rightarrow 256 \rightarrow 138$, \emph{dropout} 0,4} \\
      \hline
      Perte & \multicolumn{2}{c|}{ASL, $\gamma_- = 2$, $\gamma_+ = 0$, $c = 0{,}05$} \\
      \hline
      Optimiseur & \multicolumn{2}{c|}{AdamW, $\eta = 10^{-3}$, pénalité 0,05} \\
      \hline
      Taux de l'encodeur $\eta_{\mathrm{enc}}$ & $10^{-4}$ & $7 \cdot 10^{-5}$ \\
      \hline
      Ordonnanceur & \multicolumn{2}{c|}{décroissance cosinus jusqu'à $10^{-6}$} \\
      \hline
      Époques & 120 & 65 \\
      \hline
      Lots & \multicolumn{2}{c|}{8 molécules, 8 lots accumulés (64)} \\
      \hline
      Augmentation & \multicolumn{2}{c|}{tirage de conformation, bruit $s \in [0{,}01 ; 0{,}04]$~\AA} \\
      \hline
      Sélection du modèle & F1 de validation & AUROC de validation \\
      \hline
      Découpage & stratifié 80~/~20 & test 20~\% + 5 plis \\
      \hline
   \end{tabular}
   \end{adjustbox}
   \caption{Hyperparamètres de la meilleure expérience et du protocole d'évaluation}
   \label{tab:Hyperparametres}
\end{table}

%--------------------------------------------------------------------------------

\section{Résultats des expérimentations}
\label{sec:Resultats}

\subsection{Expériences successives}

Le tableau~\ref{tab:Experiences} présente les résultats des dix-sept expériences, extraits des journaux TensorBoard.
Les figures~\ref{fig:CourbesF1} et~\ref{fig:CourbesPerte} montrent les courbes d'apprentissage de la meilleure expérience.

\begin{table*}
   \centering
   \begin{adjustbox}{max width=\linewidth}
   \begin{tabular}{|c|l||c|c|c|}
      \hline
      N\textsuperscript{o} & Modification & F1 (0,5) & F1 (seuil opt.) & AUROC \\
      \hline
      \hline
      \multicolumn{5}{|l|}{\emph{Représentations Uni-Mol figées, découpage aléatoire 80/10/10}} \\
      \hline
      1 & Première version~: décodeur MLP, entropie croisée & 0,205 & -- & 0,839 \\
      2 & Entropie croisée pondérée & 0,258 & -- & 0,834 \\
      3 & Ordonnanceur à plateau & 0,262 & -- & 0,834 \\
      4 & AdamW, pénalité 0,01 & 0,262 & -- & 0,835 \\
      5 & AdamW, pénalité 0,05 & 0,261 & -- & 0,835 \\
      6 & \emph{Dropout} 0,5 & 0,261 & -- & 0,837 \\
      7 & Décroissance cosinus & 0,273 & -- & 0,837 \\
      \hline
      \multicolumn{5}{|l|}{\emph{Conformations 3D, encodeur gelé, découpage stratifié 80/20}} \\
      \hline
      8 & Conformations 3D et augmentation & 0,215 & -- & 0,847 \\
      9 & Seuil global optimisé & 0,209 & 0,232 & 0,844 \\
      10 & Découpage par squelettes & 0,196 & 0,209 & 0,844 \\
      11 & Sans augmentation dynamique & 0,211 & 0,234 & 0,844 \\
      12 & Sans ordonnanceur & 0,204 & 0,225 & 0,839 \\
      13 & Décodeur à 3 couches & 0,201 & 0,212 & 0,841 \\
      14 & Perte asymétrique, 300 époques & 0,249 & 0,272 & 0,857 \\
      15 & Perte asymétrique, 1\,000 époques & 0,257 & 0,288 & 0,859 \\
      \hline
      \multicolumn{5}{|l|}{\emph{Encodage corrigé}} \\
      \hline
      16 & Corrections de l'encodage, ASL réglée & 0,313 & 0,313 & 0,864 \\
      17 & Augmentation et dégel des 2 dernières couches & \textbf{0,354} & \textbf{0,354} & \textbf{0,872} \\
      \hline
   \end{tabular}
   \end{adjustbox}
   \caption{Résultats des expériences sur la validation (F1 et AUROC macro, meilleure valeur au fil des époques)}
   \label{tab:Experiences}
\end{table*}

Pour les deux dernières expériences, les métriques complémentaires sont les suivantes~:
\begin{itemize}
   \item expérience 16~: F1 macro avec seuils par descripteur de 0,374, F1 micro de 0,459~;
   \item expérience 17~: F1 macro avec seuils par descripteur de 0,400, F1 micro de 0,481.
\end{itemize}
Dans l'expérience 17, la meilleure AUROC (0,872) est atteinte à la 22\textsuperscript{e} époque et le meilleur F1 (0,354) à la 58\textsuperscript{e}~; la perte de validation est minimale dès la 15\textsuperscript{e} époque, puis remonte.

\begin{figure}
   \centering
   \placeholder{image: courbes TensorBoard de l'expérience 17 -- F1 macro et AUROC macro de validation en fonction des époques}
   \caption{Évolution du F1 et de l'AUROC de validation (expérience 17)}
   \label{fig:CourbesF1}
\end{figure}

\begin{figure}
   \centering
   \placeholder{image: courbes TensorBoard de l'expérience 17 -- perte d'entraînement et perte de validation, minimum de la validation vers l'époque 15}
   \caption{Pertes d'entraînement et de validation (expérience 17)}
   \label{fig:CourbesPerte}
\end{figure}

\begin{figure}
   \centering
   \placeholder{schema: progression du F1 macro et de l'AUROC au fil des 17 expériences, avec les trois phases (représentations figées, conformations 3D, encodage corrigé)}
   \caption{Progression des performances au fil des expériences}
   \label{fig:Progression}
\end{figure}

\subsection{Protocole d'évaluation de référence}

Le protocole d'évaluation de la section~\ref{sec:Formalisation} est implémenté et testé sur des découpages réduits, mais n'a pas encore été exécuté sur le jeu complet.
Le tableau~\ref{tab:Benchmark} est prévu pour accueillir ses résultats, à côté de ceux de la littérature.

\begin{table}
   \centering
   \begin{adjustbox}{max width=\linewidth}
   \begin{tabular}{|l||c|c|c|c|c|}
      \hline
      Modèle & AUROC & AUPRC & F1 & Précision & Rappel \\
      \hline
      \hline
      OpenPOM \cite{Han-2026} & 0,888 & 0,375 & 0,375 & 0,406 & 0,405 \\
      \hline
      Uni-Mol2 ajusté \cite{Han-2026} & 0,899 & 0,408 & 0,399 & 0,429 & 0,429 \\
      \hline
      Notre modèle (Uni-Mol v1) & \multicolumn{5}{c|}{\acompleter{résultats et IC 95~\% du protocole}} \\
      \hline
   \end{tabular}
   \end{adjustbox}
   \caption{Comparaison sur le test de GS-LF selon le protocole de \cite{Han-2026}}
   \label{tab:Benchmark}
\end{table}

%--------------------------------------------------------------------------------

\section{Analyse des résultats}
\label{sec:Analyse}

\subsection{Effet des choix successifs}

\paragraph{Représentations figées (expériences 1 à 7).}
Sans pondération, le modèle prédit très peu de descripteurs et obtient un F1 faible (0,205).
La pondération des positifs a le plus fort effet (0,258)~; la régularisation (AdamW, \emph{dropout}) et l'ordonnanceur cosinus apportent ensuite de petits gains, jusqu'à 0,273.
Le surapprentissage reste marqué~: la perte de validation est minimale dès les premières époques.
Une représentation figée, calculée une fois par molécule, ne permet pas d'augmenter les données.

\paragraph{Conformations 3D (expériences 8 à 15).}
Le passage aux conformations et à l'augmentation élimine le surapprentissage et améliore l'AUROC (0,847), mais le F1 au seuil de 0,5 baisse.
Ces deux mesures ne sont toutefois pas directement comparables, la validation ayant changé (découpage 80/20 stratifié au lieu de 80/10/10 aléatoire).
Le seuil de 0,5 est par ailleurs inadapté à des descripteurs rares~: l'optimiser apporte environ deux points de F1.
Le découpage par squelettes fait perdre environ 2 points, ce qui indique une part de mémorisation des familles chimiques.
L'augmentation dynamique, l'ordonnanceur et un décodeur plus profond n'apportent pas de gain mesurable dans cette configuration.
La perte asymétrique apporte le gain le plus net de cette phase~: l'AUROC passe à 0,857 et le F1 à seuil optimisé à 0,272, puis 0,288 avec un entraînement plus long.

\paragraph{Encodage corrigé (expériences 16 et 17).}
La correction de l'encodage, combinée à un réglage de la perte asymétrique, porte le F1 à 0,313 en 66 époques, contre 0,288 en 800 époques auparavant.
Ces deux changements ayant été faits ensemble, leurs effets respectifs ne peuvent pas être séparés~; la rapidité de convergence suggère toutefois que l'encodage incorrect pénalisait fortement les expériences précédentes.
Avec l'encodeur gelé, le F1 plafonne~: c'est le dégel des deux dernières couches, avec l'augmentation, qui le porte à 0,354 et l'AUROC à 0,872.
En contrepartie, les 6,3 millions de paramètres dégelés provoquent un surapprentissage précoce~: la perte de validation remonte après la 15\textsuperscript{e} époque, alors que le F1 continue de progresser jusqu'à la 58\textsuperscript{e}.

\subsection{Comparaison avec la littérature}

Le F1 macro de 0,354 de l'expérience 17 est proche de celui de POM (0,360) et d'OpenPOM (0,375).
La comparaison n'est cependant pas valide en l'état, pour quatre raisons~:
\begin{enumerate}
   \item nos métriques sont mesurées sur la validation, qui sert aussi à choisir l'époque, et non sur un test tenu à l'écart~;
   \item nous rapportons la meilleure valeur au fil des époques, ce qui est optimiste~;
   \item le F1 avec seuils par descripteur (0,400) est encore plus optimiste, les seuils étant choisis sur les données mêmes où il est mesuré~; dans le protocole de référence, ils sont choisis sur les prédictions hors pli puis appliqués au test~;
   \item les modèles de référence sont des ensembles (50 modèles pour Uni-Mol2), un unique découpage ne permet pas d'estimer l'incertitude, et notre validation (environ 1\,000 molécules) n'est pas celle de leur test.
\end{enumerate}

À l'inverse, nos résultats sont obtenus avec Uni-Mol v1 et sans recherche d'hyperparamètres~; ils pourraient donc progresser.
Enfin, il faut distinguer F1 macro et F1 micro~: le F1 micro, dominé par les descripteurs fréquents, atteint 0,481, mais c'est le F1 macro que publie la littérature.
Seule l'exécution du protocole (tableau~\ref{tab:Benchmark}) permettra de conclure.

\subsection{Limites des expériences}

Chaque expérience ne correspond qu'à une exécution, avec une graine aléatoire fixe~: les écarts de l'ordre du point de F1 ne sont pas significatifs.
Plusieurs expériences modifient plus d'un paramètre à la fois.
Les expériences 8 à 15 ont été menées avec un encodage incorrect, et leurs conclusions (inutilité de l'augmentation dynamique, du décodeur profond) mériteraient d'être revérifiées avec l'encodage corrigé.

%--------------------------------------------------------------------------------

\section{Conclusion}
\label{sec:ConclusionExperimentations}

La chaîne développée est fonctionnelle et a permis de doubler presque le F1 macro de validation, de 0,205 à 0,354, et de porter l'AUROC à 0,872.
Les choix les plus déterminants ont été le traitement du déséquilibre (pondération puis perte asymétrique), la conformité de l'encodage des entrées à celui du pré-entraînement et l'ajustement fin des dernières couches de l'encodeur.
Ce dernier introduit un surapprentissage précoce, qu'il faudra maîtriser.

Le principal enseignement méthodologique est que nos premières mesures n'étaient pas comparables à la littérature.
Le protocole de référence est désormais implémenté~; son exécution est la première étape de la phase~II.
Pour un repreneur, nous recommandons de vérifier l'encodage des entrées d'un modèle pré-entraîné avant toute expérience, et de fixer le protocole d'évaluation dès le début du projet.

\chapter{Conclusion}
\label{chap:Conclusion}

À l'issue de cette première phase, nous résumons le travail effectué, en tirons les principaux enseignements et présentons les perspectives de la phase~II.

%--------------------------------------------------------------------------------

\section{Résumé du travail effectué}
\label{sec:Resume}

Nous avons étudié l'apport d'un transformeur de graphes 3D pré-entraîné, Uni-Mol, à la prédiction multi-label des 138 descripteurs olfactifs du jeu GS-LF.
L'étude bibliographique a situé ce choix par rapport aux GNN de référence, POM et OpenPOM, limités par la myopie de la propagation de messages et l'absence de géométrie, et nous a fait découvrir un travail récent appliquant Uni-Mol2 à la même tâche, dont nous avons repris le protocole d'évaluation.

Nous avons développé une chaîne d'apprentissage complète~: génération de conformations 3D, encodage au format Uni-Mol, augmentation géométrique, décodeur MLP, perte asymétrique et dégel partiel de l'encodeur.
Dix-sept expériences ont porté le F1 macro de validation de 0,205 à 0,354 et l'AUROC macro à 0,872.
Ces valeurs sont du même ordre que celles des modèles de référence (F1 de 0,375 pour OpenPOM, 0,399 pour Uni-Mol2 ajusté), mais ne leur sont pas encore comparables~: elles sont mesurées sur la validation, sans test tenu à l'écart.
Le protocole de référence, implémenté en fin de phase, permettra cette comparaison.

%--------------------------------------------------------------------------------

\section{Enseignements}
\label{sec:Enseignements}

\paragraph{Le protocole d'évaluation fait partie du résultat.}
Un F1 n'a de sens qu'avec le découpage, le choix des seuils et la sélection du modèle qui l'accompagnent.
Nous avons longtemps optimisé une mesure optimiste~; fixer le protocole de la littérature dès le début aurait permis des comparaisons fiables tout au long du projet.
De même, F1 micro et F1 macro mesurent des choses différentes~: le premier, dominé par les descripteurs fréquents, est bien plus élevé que le second.

\paragraph{Un modèle pré-entraîné exige un encodage exact.}
Les écarts d'encodage que nous avons corrigés ne produisaient aucune erreur, seulement des performances dégradées.
La correction a apporté l'un des gains les plus importants du projet~: vérifier la conformité des entrées au pré-entraînement est une étape à part entière.

\paragraph{Le déséquilibre se traite à plusieurs niveaux.}
La pondération des positifs, puis la perte asymétrique, ont apporté des gains nets, et le choix des seuils de décision en apporte d'autres.
Ces leviers agissent à des niveaux différents (apprentissage, décision) et se combinent.

\paragraph{Peu de données, beaucoup de paramètres.}
Le dégel de l'encodeur améliore les performances mais provoque un surapprentissage rapide~: sur 4\,000 molécules, chaque million de paramètres libérés doit être justifié.

\paragraph{Sur le plan de l'organisation.}
Le suivi systématique des expériences (configuration, courbes et commentaire pour chacune) a rendu ce rapport possible et permet de revenir sur des conclusions antérieures.
À l'inverse, modifier plusieurs paramètres à la fois a empêché d'attribuer certains gains.
\acompleter{enseignements personnels de chaque auteur}

%--------------------------------------------------------------------------------

\section{Perspectives de recherche}
\label{sec:Perspectives}

\paragraph{Évaluer selon le protocole de référence.}
La première étape de la phase~II est d'exécuter le protocole sur le jeu complet, afin de situer réellement notre modèle par rapport à POM, OpenPOM et Uni-Mol2.
Il faudra ensuite maîtriser le surapprentissage lié au dégel (taux d'apprentissage plus faible pour l'encodeur, sélection du modèle sur l'AUROC, dégel progressif), revérifier les conclusions des expériences antérieures à la correction de l'encodage, et mener une recherche d'hyperparamètres.

\paragraph{Passer à Uni-Mol2.}
Le sujet désigne Uni-Mol2 comme modèle cible.
Sa représentation de paires, qui combine géométrie et topologie, et sa taille supérieure devraient améliorer les performances, comme le suggère \cite{Han-2026}, au prix d'un coût de calcul plus élevé.

\paragraph{Modéliser les dépendances entre descripteurs.}
Le décodeur traite les 138 descripteurs comme indépendants~: chaque ligne $\mathbf{w}_j$ de la dernière matrice $\mathbf{W} \in \mathbb{R}^{138 \times 256}$ est apprise uniquement à partir des molécules qui portent ou non le descripteur $j$.
Les descripteurs rares (31 exemples pour le plus rare) apprennent donc mal, alors que les confusions du modèle suivent les co-occurrences entre descripteurs \cite{Han-2026}.
Nous proposons de contraindre la matrice sous la forme
\begin{equation}
   \label{eq:Factorisation}
   \mathbf{W} = \mathbf{E} \, \mathbf{P},
   \qquad
   \mathbf{E} \in \mathbb{R}^{138 \times d},
   \quad
   \mathbf{P} \in \mathbb{R}^{d \times 256},
   \quad
   d < 138
\end{equation}
où $\mathbf{E}$ est figée et décrit chaque descripteur, par exemple par une décomposition en valeurs singulières de la matrice de co-occurrence calculée sur l'ensemble d'entraînement, et où seule $\mathbf{P}$ est apprise.
Chaque $\mathbf{w}_j = \mathbf{e}_j \mathbf{P}$ est alors une combinaison de $d$ vecteurs partagés~: deux descripteurs aux profils de co-occurrence proches ont nécessairement des vecteurs proches, et les descripteurs rares bénéficient des exemples de leurs voisins.
La condition $d < 138$ est essentielle~: si $d = 138$ et $\mathbf{E}$ est inversible, tout $\mathbf{W}$ reste atteignable et la contrainte disparaît.
La dimension $d$ règle donc le degré de partage imposé entre descripteurs.
Cette approche est un premier pas vers une représentation conjointe des molécules et des descripteurs, qui permettrait aussi de prédire des descripteurs absents de l'entraînement.

\paragraph{Agréger des données, notamment stéréochimiques.}
GS-LF est déjà la fusion des catalogues GoodScents et Leffingwell~; récupérer leurs versions brutes, par exemple via la plateforme Pyrfume, apporterait peu de nouvelles molécules.
Leur intérêt est ailleurs~: GS-LF a perdu la stéréochimie lors de sa constitution (les SMILES sont non stéréospécifiques), et les sources d'origine pourraient la contenir, ainsi qu'un vocabulaire plus riche.
D'autres bases (Dravnieks, Keller, Arctander) apporteraient de nouvelles molécules, avec leurs propres vocabulaires.
L'agrégation impose d'harmoniser ces vocabulaires, de résoudre les annotations contradictoires et, surtout, d'exclure toute molécule du test de GS-LF pour préserver la comparabilité.
La factorisation~\eqref{eq:Factorisation} s'y prête naturellement~: chaque source peut avoir ses propres lignes dans $\mathbf{E}$ tout en partageant $\mathbf{P}$.

L'objectif est de donner un réel avantage à l'encodeur 3D par rapport à un encodeur topologique.
Mais une limite doit être levée~: Uni-Mol ne voit que des distances, qui sont identiques pour deux énantiomères.
Les données chirales n'aideront donc que pour les diastéréoisomères, sauf à fournir explicitement la configuration des centres stéréogènes au modèle, par exemple comme caractéristique de chaque atome.
Un test simple, comparer les sorties du modèle pour les deux énantiomères de la carvone, permettra de vérifier ce point avant d'engager l'agrégation.

\paragraph{Expliquer les prédictions.}
Le sujet demande d'expliquer le fonctionnement du modèle sur quelques molécules.
Les cartes d'attention d'Uni-Mol et des méthodes d'attribution par atome permettraient de vérifier si les prédictions s'appuient sur des groupes fonctionnels connus pour leur rôle olfactif (esters pour le fruité, composés soufrés, etc.).

\paragraph{Compléter l'état de l'art.}
Les travaux GraphNOSE \cite{Sharma-2026} et sur les représentations à double graphe \cite{Zhang-2026}, identifiés en fin de phase, seront étudiés et comparés à notre approche.


%--------------------------------------------------------------------------------

\bibliography{rapport}

\listoffigures{}

\listoftables{}

\listofalgorithms{}

\appendix

\chapter{Fiches de lecture}
\label{ann:FichesLecture}

Le tableau~\ref{tab:FichesLecture} récapitule les travaux étudiés et leur lien avec le projet~; chacun fait ensuite l'objet d'une fiche.

\begin{table}[H]
   \centering
   \begin{adjustbox}{max width=\linewidth}
   \begin{tabular}{|l|c|p{0.5\linewidth}|}
      \hline
      Référence & Année & Apport pour le projet \\
      \hline
      \hline
      Sanchez-Lengeling et al. \cite{SanchezLengeling-2019} & 2019 & Formulation multi-label de la tâche, premier GNN \\
      \hline
      Lee et al. \cite{Lee-2023} & 2023 & Jeu GS-LF, modèle de référence POM \\
      \hline
      Zhou et al. \cite{Zhou-2023} & 2023 & Encodeur Uni-Mol utilisé \\
      \hline
      Ji et al. \cite{Ji-2024} & 2024 & Uni-Mol2, modèle cible du sujet \\
      \hline
      Han et al. \cite{Han-2026} & 2026 & Uni-Mol2 en olfaction, protocole d'évaluation repris \\
      \hline
      Ridnik et al. \cite{Ridnik-2021} & 2021 & Perte asymétrique utilisée \\
      \hline
   \end{tabular}
   \end{adjustbox}
   \caption{Récapitulatif des fiches de lecture}
   \label{tab:FichesLecture}
\end{table}

\paragraph{\emph{Machine learning for scent: learning generalizable perceptual representations of small molecules}.}
Premier travail d'apprentissage profond à grande échelle sur la prédiction des odeurs \cite{SanchezLengeling-2019}.

\subparagraph{Résumé.}
Les auteurs constituent un jeu d'environ 5\,000 molécules annotées par 138 descripteurs, issu des catalogues GoodScents et Leffingwell, et entraînent un GNN par propagation de messages sur le graphe moléculaire.
Le GNN surpasse des modèles fondés sur des empreintes moléculaires et des descripteurs chimiques.
Sa représentation interne regroupe les molécules selon leurs odeurs et se transfère à d'autres tâches olfactives.

\subparagraph{Analyse.}
Ce travail fonde la formulation que nous reprenons (multi-label, 138 descripteurs, métriques macro).
Il utilise un modèle 2D appris à partir de zéro, ce qui laisse ouverte la question de l'apport de la géométrie et du pré-entraînement, au cœur de notre sujet.

\paragraph{\emph{A principal odor map unifies diverse tasks in olfactory perception}.}
Prolongement du travail précédent, publié dans \emph{Science} \cite{Lee-2023}.

\subparagraph{Résumé.}
Le modèle POM, un GNN entraîné sur la fusion nettoyée GS-LF, définit une carte des odeurs.
Une validation prospective sur environ 400 molécules nouvelles, évaluées par un panel entraîné, montre que ses prédictions sont, pour une majorité des molécules, plus proches du consensus du panel que celles d'un panéliste médian.
La carte se transfère à d'autres tâches (intensité, similarité, etc.).

\subparagraph{Analyse.}
POM est le modèle de référence de la tâche et son protocole (test tenu à l'écart, validation croisée, ensemble) est repris par les travaux ultérieurs.
Son code n'étant pas publié, la comparaison reproductible passe par OpenPOM.
Pour notre projet, il fournit les données et les performances à atteindre.

\paragraph{\emph{Uni-Mol: a universal 3D molecular representation learning framework}.}
Proposition d'un transformeur moléculaire 3D pré-entraîné \cite{Zhou-2023}.

\subparagraph{Résumé.}
Uni-Mol traite les atomes comme des tokens et encode leurs distances par noyaux gaussiens dans une représentation de paires qui biaise l'attention.
Il est pré-entraîné sur plus de 200 millions de conformations de 19 millions de molécules, en retrouvant des atomes masqués et en débruitant des coordonnées.
Ajusté sur des tâches de prédiction de propriétés moléculaires, il obtient de très bons résultats.

\subparagraph{Analyse.}
C'est l'encodeur utilisé dans notre travail, via la bibliothèque \texttt{unimol\_tools}.
Sa représentation, fondée sur les distances, est invariante par réflexion~: elle ne distingue pas les énantiomères, limite importante pour l'olfaction.
L'encodage des entrées doit reproduire exactement celui du pré-entraînement.

\paragraph{\emph{Exploring molecular pretraining model at scale}.}
Article principal désigné par le sujet \cite{Ji-2024}.

\subparagraph{Résumé.}
Les auteurs constituent un jeu d'environ 884 millions de molécules et entraînent Uni-Mol2, un transformeur à deux pistes (atomes et paires), de 84 millions à 1,1 milliard de paramètres.
Les paires combinent distances, types de liaisons et plus courts chemins dans le graphe.
Ils établissent une loi d'échelle reliant la perte de validation à la taille du modèle, des données et du calcul, et montrent que les performances en aval progressent avec la taille.

\subparagraph{Analyse.}
Uni-Mol2 est la cible de la phase~II.
L'ajout d'informations topologiques aux distances devrait aider, mais la stéréochimie des énantiomères reste a priori hors de portée.
Le coût de calcul impose de commencer par la plus petite version (84M), comme l'a fait \cite{Han-2026}.

\paragraph{\emph{A general-purpose molecular foundation model transfers across diverse olfactory tasks}.}
Application d'Uni-Mol2 à l'olfaction \cite{Han-2026}.

\subparagraph{Résumé.}
Uni-Mol2 (84M) est ajusté une fois sur GS-LF avec une perte focale, puis évalué sans nouvel entraînement sur plusieurs tâches olfactives.
Sur le test de GS-LF, il dépasse POM et OpenPOM sur les cinq métriques macro, avec des écarts modestes.
Il distingue les énantiomères, contrairement à OpenPOM, mais sans savoir lequel porte quelle odeur.
L'analyse des erreurs montre que les confusions suivent les co-occurrences entre descripteurs.

\subparagraph{Analyse.}
C'est le travail le plus proche du nôtre~: nous en reprenons le protocole d'évaluation pour comparer nos résultats.
Plusieurs détails d'entraînement manquent pour le reproduire exactement.
Son analyse des confusions inspire notre proposition de factorisation du décodeur.

\paragraph{\emph{Asymmetric loss for multi-label classification}.}
Proposition d'une fonction de perte adaptée au multi-label déséquilibré \cite{Ridnik-2021}.

\subparagraph{Résumé.}
En multi-label, chaque exemple a beaucoup plus de labels négatifs que positifs.
La perte asymétrique applique une focalisation différente aux positifs et aux négatifs, et décale la probabilité des négatifs d'une marge qui annule la perte des plus faciles et atténue l'effet des étiquettes erronées.
Elle améliore l'état de l'art sur plusieurs bancs d'essai de classification d'images multi-label.

\subparagraph{Analyse.}
Notre problème présente exactement ce déséquilibre (environ 4,9 descripteurs positifs sur 138 par molécule).
La perte asymétrique a apporté l'un des gains les plus nets de nos expériences.

%--------------------------------------------------------------------------------

\chapter{Planification}
\label{ann:Planification}

La figure~\ref{fig:PlanningPrevisionnel} présente le planning élaboré au début du projet.

\begin{figure*}
   \centering
   \placeholder{schema: diagramme de Gantt prévisionnel de la phase I}
   \caption{Planification prévisionnelle}
   \label{fig:PlanningPrevisionnel}
\end{figure*}

La figure~\ref{fig:PlanningEffectif} présente le planning effectif, relevé au fur et à mesure de l'avancement.

\begin{figure*}
   \centering
   \placeholder{schema: diagramme de Gantt effectif -- mise en place du projet (13--15 oct.), représentations figées et premières expériences (15--17 oct.), conformations 3D (18--21 oct.), perte asymétrique (25 oct.--3 nov.), corrections de l'encodage et dégel (4--6 nov.), suivi des expériences et schémas (16 nov.), validation croisée et protocole d'évaluation (17--24 nov.)}
   \caption{Planning effectif}
   \label{fig:PlanningEffectif}
\end{figure*}

\acompleter{discussion des écarts entre les deux plannings et leçons sur la gestion du projet}

%--------------------------------------------------------------------------------

\chapter{Fiches de suivi}
\label{ann:FichesSuivi}

% Les temps de travail (\tempstravailA pour Tom Daniaud, \tempstravailB pour Florian Petiot)
% sont à compléter : valeurs à 0 par défaut.
% Le travail effectué est reconstitué à partir de l'historique Git.

\begin{fichesuivi}{}{12 octobre 2026}
   \tempstravailA{0}{0}
   \tempstravailB{0}{0}

   \begin{travaileffectue}
      \begin{itemize}
         \item \acompleter{étude préalable et bibliographique antérieure au développement}
      \end{itemize}
   \end{travaileffectue}

   \begin{travailnoneffectue}
   \end{travailnoneffectue}

   \begin{echange}
      \begin{itemize}
         \item \acompleter{réunions avec les encadrants}
      \end{itemize}
   \end{echange}

   \begin{planification}
   \end{planification}
\end{fichesuivi}

\begin{fichesuivi}{13 octobre 2026}{17 octobre 2026}
   \tempstravailA{0}{0}
   \tempstravailB{0}{0}

   \begin{travaileffectue}
      \begin{itemize}
         \item mise en place du dépôt, de la configuration et du suivi TensorBoard~;
         \item intégration du jeu GS-LF et calcul des représentations Uni-Mol figées~;
         \item décodeur de classification, métriques et premières expériences (1 à 7)~;
         \item analyse des données et pondération des descripteurs~;
         \item création du journal d'expériences.
      \end{itemize}
   \end{travaileffectue}

   \begin{travailnoneffectue}
   \end{travailnoneffectue}

   \begin{echange}
   \end{echange}

   \begin{planification}
      \begin{itemize}
         \item passer à des conformations 3D pour permettre l'augmentation de données.
      \end{itemize}
   \end{planification}
\end{fichesuivi}

\begin{fichesuivi}{18 octobre 2026}{24 octobre 2026}
   \tempstravailA{0}{0}
   \tempstravailB{0}{0}

   \begin{travaileffectue}
      \begin{itemize}
         \item génération hors ligne des conformations 3D et intégration d'Uni-Mol comme encodeur~;
         \item optimisation du seuil de décision~;
         \item découpage par squelettes de Bemis-Murcko~;
         \item prise en charge des GPU Intel Arc~;
         \item expériences 8 à 13.
      \end{itemize}
   \end{travaileffectue}

   \begin{travailnoneffectue}
   \end{travailnoneffectue}

   \begin{echange}
   \end{echange}

   \begin{planification}
   \end{planification}
\end{fichesuivi}

\begin{fichesuivi}{25 octobre 2026}{31 octobre 2026}
   \tempstravailA{0}{0}
   \tempstravailB{0}{0}

   \begin{travaileffectue}
      \begin{itemize}
         \item ajout du F1 micro~;
         \item réorganisation des pertes et métriques, perte asymétrique~;
         \item expérience 14.
      \end{itemize}
   \end{travaileffectue}

   \begin{travailnoneffectue}
   \end{travailnoneffectue}

   \begin{echange}
   \end{echange}

   \begin{planification}
   \end{planification}
\end{fichesuivi}

\begin{fichesuivi}{1\textsuperscript{er} novembre 2026}{7 novembre 2026}
   \tempstravailA{0}{0}
   \tempstravailB{0}{0}

   \begin{travaileffectue}
      \begin{itemize}
         \item décodeur à plusieurs couches~;
         \item correction de l'encodage Uni-Mol (dictionnaire, distances de \texttt{[CLS]}, mode évaluation de l'encodeur gelé)~;
         \item taux d'apprentissage différencié pour l'encodeur et F1 par descripteur~;
         \item expériences 15 à 17.
      \end{itemize}
   \end{travaileffectue}

   \begin{travailnoneffectue}
   \end{travailnoneffectue}

   \begin{echange}
   \end{echange}

   \begin{planification}
   \end{planification}
\end{fichesuivi}

\begin{fichesuivi}{8 novembre 2026}{14 novembre 2026}
   \tempstravailA{0}{0}
   \tempstravailB{0}{0}

   \begin{travaileffectue}
      \begin{itemize}
         \item \acompleter{travail de la semaine}
      \end{itemize}
   \end{travaileffectue}

   \begin{travailnoneffectue}
   \end{travailnoneffectue}

   \begin{echange}
   \end{echange}

   \begin{planification}
   \end{planification}
\end{fichesuivi}

\begin{fichesuivi}{15 novembre 2026}{21 novembre 2026}
   \tempstravailA{0}{0}
   \tempstravailB{0}{0}

   \begin{travaileffectue}
      \begin{itemize}
         \item réorganisation des journaux d'expériences~;
         \item schémas de la chaîne et de l'architecture (Excalidraw)~;
         \item validation croisée à cinq plis.
      \end{itemize}
   \end{travaileffectue}

   \begin{travailnoneffectue}
   \end{travailnoneffectue}

   \begin{echange}
   \end{echange}

   \begin{planification}
   \end{planification}
\end{fichesuivi}

\begin{fichesuivi}{22 novembre 2026}{24 novembre 2026}
   \tempstravailA{0}{0}
   \tempstravailB{0}{0}

   \begin{travaileffectue}
      \begin{itemize}
         \item étude de \cite{Han-2026} et comparaison avec la littérature~;
         \item découpage de référence et implémentation du protocole d'évaluation~;
         \item rédaction du rapport de phase~I.
      \end{itemize}
   \end{travaileffectue}

   \begin{travailnoneffectue}
      \begin{itemize}
         \item exécution complète du protocole d'évaluation, reportée à la phase~II.
      \end{itemize}
   \end{travailnoneffectue}

   \begin{echange}
      \begin{itemize}
         \item pistes pour la phase~II~: modélisation des dépendances entre descripteurs, agrégation de données stéréochimiques.
      \end{itemize}
   \end{echange}

   \begin{planification}
      \begin{itemize}
         \item exécuter le protocole d'évaluation~;
         \item passer à Uni-Mol2.
      \end{itemize}
   \end{planification}
\end{fichesuivi}

\printweeksummary

%--------------------------------------------------------------------------------

\chapter{Auto-contrôle et auto-évaluation}
\label{ann:AutoEvaluation}

La figure~\ref{fig:AutoEvaluationTravailIntermediaire} énumère les points à contrôler dans les trois composantes du travail (rapport, présentation orale, travail de fond) et notre niveau de satisfaction à l'issue de la phase~I.

\begin{figure*}
   \centering
      \ifscreen
         \rotatebox{90}{\includegraphics[width=0.9\textheight]{Images/Grille-Evaluation-PRD1}}
      \else
         \includegraphics[width=0.9\textwidth]{Images/Grille-Evaluation-PRD1}
      \fi
   \caption{Points à contrôler à l'issue de la phase I}
   \label{fig:AutoEvaluationTravailIntermediaire}
\end{figure*}

\acompleter{commentaire des points de satisfaction et d'insatisfaction}


\end{document}
